\documentclass[journal,twoside,web]{ieeecolor}

\include{Config/config_ieee}
\include{Config/config_des}

% abbrs for this article only
\newcommand{\sFourC}{\ensuremath{\text{S}^4\text{C}}}

\markboth{\hskip25pc IEEE TRANSACTIONS AND JOURNALS TEMPLATE}
{Author \MakeLowercase{\textit{et al.}}: Title}
\begin{document}
\title{Synthesis of Cognitive Non-Blocking Supervisors for Partially-Observed Discrete-Event Systems}
\author{Yiwei Zheng, Bohan Cui, and Xiang Yin
\thanks{TODO This paragraph of the first footnote will contain the date on 
which you submitted your article for review. It will also contain support 
information, including sponsor and financial support acknowledgment. For 
example, ``This work was supported in part by the U.S. Department of 
Commerce under Grant 123456.'' }
\thanks{TODO First A. Author is with the National Institute of Standards and 
Technology, Boulder, CO 80305 USA (e-mail: author@boulder.nist.gov). }
\thanks{TODO Second B. Author Jr. was with Rice University, Houston, TX 77005 USA. He is 
now with the Department of Physics, Colorado State University, Fort Collins, 
CO 80523 USA (e-mail: author@lamar.colostate.edu).}}

\maketitle

\begin{abstract}
This article investigates the problem of synthesizing cognitive non-blocking supervisors for partially observed discrete event systems (DESs).
The notion of cognitive non-blockingness requires that for each state of the DES, the system can eventually certify its arrival at a target state based solely on partial observations.
The synthesis problem is therefore to construct a supervisor under which the closed-loop system satisfies this observation-based progress requirement.
%
We show that the non-blocking all-inclusive controller (AIC) proposed in \cite{yin2016synthesis} is insufficient for this purpose, and introduce a new structure, named the bipartite cognitive controller (BCC), to capture the microscopic evolution within the AIC.
We prove that the event sequences generated by a BCC can be directly projected as the language of the controlled DES, thereby establishing the equivalence between the BCC-induced language and the controlled DES language.
Moreover, we show that cognitive non-blockingness can be guaranteed by requiring the Y-states to be cognitively live and the Z-states to be potentially feasible.
The resulting BCC can then be projected back to a supervisor for cognitive non-blocking synthesis of the original DES.
\end{abstract}

\begin{IEEEkeywords}
Discrete-event systems (DES), partial observation, non-blocking, supervisory control. %maximal permissiveness.
%\underline{https://www.ieee.org/publications/services/thesaurus-access-page.com}
\end{IEEEkeywords}

\section{Introduction}
\label{sec:intro}
\IEEEPARstart{A}{s} modern computation and networking technologies continue to advance, discrete event systems (DESs) have become increasingly important in various aspects of modern life and industrial production \cite{yin2024formal, fokkink2024offline}.
Owing to their capability to model complex event-driven behaviors, DESs have been widely applied in areas such as automated manufacturing systems, communication networks, transportation systems, and robotic coordination.
%
Non-blockingness is one of the most fundamental properties of DESs, since it guarantees that the system can always complete its prescribed tasks without being trapped before task completion \cite{cassandras2021introduction}.
In general, a DES is said to be non-blocking if every reachable state can still evolve to a marked (or target) state through some feasible event sequence.
Equivalently, every generated behavior of the system admits a continuation that reaches a marked behavior.
Therefore, non-blockingness ensures that the system can continuously make progress toward task completion rather than entering an irreversible blocking state, for example, a deadlock or a livelock state.

\begin{figure}[h!]
    \centering
    \includegraphics[width=0.65\linewidth]{Figs/motivated_example_overall.pdf}
    \caption{Motivated example of a robot deployed in an unknown environment.}
    \label{fig:moti_ex:1}
\end{figure}

%
% TODO
% 要不要把enforce这个词引入?
The supervisory control framework introduced by Ramadge and Wonham \cite{ramadge1987supervisory} provides a systematic approach to enforcing desired behaviors of DESs by enabling or disabling controllable events.
In the classical non-blocking synthesis problem, the objective is to construct a supervisor such that every behavior generated by the closed-loop system can be extended to a marked behavior.
%
However, under partial observation, classical non-blockingness may be insufficient from the supervisor's perspective.
Since the supervisor cannot observe the exact current state of the plant, it can only maintain the set of states consistent with the observed event sequence.
Thus, even if the actual state can eventually reach a marked state, the supervisor may not be able to certify from its observations that a marked state has been reached.
%
Let us consider the motivated example of a robotic system deployed in a partially observable environment, which is illustrated in Fig.~\ref{fig:moti_ex:1}. 
The robot knows the environment model, but its exact location is uncertain due to limited sensing.
Suppose that the robot is allowed to transmit data only in safe regions, which are regarded as the target regions in this task.
If a safe target location is observationally indistinguishable from an unsafe location, then reaching a safe region in the actual system is not enough.
Before transmitting data, the robot must be able to certify from its observation that all possible current locations are safe target locations.

To address the above limitation, this article develops a new notion named \emph{cognitive non-blocking}, which aims to guarantee not only task reachability but also the observational certifiability of task completion under partial observation.
%
In contrast to classical non-blockingness, which only requires that every closed-loop behavior can be extended to some marked behavior, cognitive non-blocking further requires that, after some admissible continuation, the supervisor can certify from its observation that the current state belongs to the marked target set.
%
Built upon the supervisory control framework, this article further studies the synthesis of partial-observation supervisors that enforce cognitive non-blockingness on the closed-loop system.
Different from classical non-blocking synthesis, the proposed problem must account for the evolution of information states, since task completion has to be certified from observations rather than merely reached by the actual plant state.

The main contributions of this article are summarized as follows.
%
\begin{enumerate}[(i)]
    \item We propose a new observation-based property, named \emph{cognitive non-blocking}, for partially observed DESs.
    Compared with classical non-blockingness studied in  \cite{yin2016synthesis,yin2018synthesis}, the proposed notion strengthens the standard reachability-based requirement by incorporating observational certification of task completion.
    Specifically, it requires not only that marked states are eventually reachable, but also that the supervisor can certify, based solely on partial observations, that the current state belongs to the marked set.

    \item We develop a new structural representation, called the     \emph{bipartite cognitive controller}, to characterize the microscopic     evolution of information states within a bipartite controller.
    Existing bipartite-controller-based approaches      \cite{yin2016uniform,zheng2026enforcement} describe the closed-loop     evolution in terms of control decisions and observations.
    In contrast, the proposed bipartite cognitive controller further refines this evolution at the microstate level, which enables the verification and synthesis of cognitive non-blocking supervisors.
    By iteratively pruning Z-states that are not potentially feasible and Y-states that are not cognitively live, while preserving the consistency of the remaining structure, we obtain a controlled structure whose induced language is a sublanguage of the plant language and whose closed-loop behavior is guaranteed to satisfy cognitive non-blockingness.
    
\end{enumerate}

\subsection{Related works}

\begin{table}[!t]
\centering
\caption{Literature review on non-blockingness and supervisory control for DESs.}
\label{tab:literature_review}
\renewcommand{\arraystretch}{1.15}
\setlength{\tabcolsep}{2.5pt}
\begin{tabular}{
@{}
L{0.18\columnwidth}
L{0.43\columnwidth}
L{0.27\columnwidth}
@{}}
\toprule
\textbf{Dimension}
& \textbf{Focus}
& \textbf{Representative} \\
&
& \textbf{references} \\
\midrule

\multirow{4}{0.18\columnwidth}{System model}
    & Finite-state DES
    & \cite{ramadge1987supervisory}, \cite{lin1988observability},
      \cite{chen1991dealing}, \cite{yin2016uniform},
      \cite{yin2016synthesis}, \cite{yin2018synthesis},
    \cite{xiang2024existence},
      \cite{zhang2024quantitatively} \\

    & Modular / compositional DES
    & \cite{malik2013compositional}, \cite{mohajerani2016framework},
      \cite{pilbrow2015algorithm}, \cite{tomas2018complexity},
      \cite{tang2024compositional}, \cite{thuijsman2024transformational},
      \cite{jan2024supervisory} \\

    & Timed / networked DES
    & \cite{basile2022supervisory}, \cite{aida2020blocking},
      \cite{hou2023distributed},
      \cite{zhang2025bounded}, \cite{hu2026supervisory} \\

    & Extended DES models
    & \cite{li2025supervisory}, \cite{zheng2026enforcement},
      \cite{zhang2024quantitatively}, \cite{liu2025approximate}, \cite{hu2026supervisory} \\
\midrule

\multirow{4}{0.18\columnwidth}{Observation setting}
    & Full-observation supervisory control
    & \cite{ramadge1987supervisory}, \cite{chen1991dealing},
      \cite{malik2013compositional}, \cite{tomas2018complexity},
      \cite{tang2024compositional} \\

    & Partial-observation supervisory control
    & \cite{lin1988observability}, \cite{yin2016uniform},
      \cite{yin2016synthesis}, \cite{yin2018synthesis},
      \cite{hu2024design},
      \cite{jan2024supervisory}, \cite{hu2026supervisory} \\

    & Information-state-based control
    & \cite{yin2016uniform}, \cite{yin2016synthesis},
    \cite{hu2024design},
      \cite{xie2023optimal}, \cite{zheng2026enforcement},
      \cite{liu2025approximate}\\

    & Observation-based property enforcement
    & \cite{shu2013enforcing}, \cite{hu2023asynchronous},
      \cite{xie2023optimal}, \cite{li2025supervisory},
      \cite{zheng2026enforcement}, \cite{hu2026supervisory} \\
\midrule

\multirow{4}{0.18\columnwidth}{Property / problem}
    & Classical non-blockingness
    & \cite{chen1991dealing}, \cite{yin2016synthesis},
      \cite{yin2018synthesis} \\

    & Modular non-blocking verification
    & \cite{malik2013compositional}, \cite{mohajerani2016framework},
      \cite{pilbrow2015algorithm}, \cite{tomas2018complexity},
      \cite{tang2024compositional}, \cite{thuijsman2024transformational} \\

    & Quantitative / bounded / timed non-blockingness
    & \cite{xiang2024existence}, \cite{zhang2024quantitatively},
      \cite{zhang2025bounded}, \cite{hou2023distributed} \\

    & Other IS-based properties
    & \cite{shu2013enforcing}, \cite{ma2022resilient},
      \cite{hu2023asynchronous}, \cite{xie2023optimal}, 
      \cite{li2025supervisory}, \cite{zheng2026enforcement}, 
      \cite{liu2025approximate}, \cite{hu2026supervisory} \\
\midrule

\multirow{3}{0.18\columnwidth}{Contribution type}
    & Verification / complexity analysis
    & \cite{chen1991dealing}, \cite{malik2013compositional},
      \cite{mohajerani2016framework}, \cite{pilbrow2015algorithm},
      \cite{tomas2018complexity}, \cite{tang2024compositional},
      \cite{thuijsman2024transformational}, \cite{liu2025approximate} \\

    & Supervisor synthesis
    & \cite{ramadge1987supervisory}, \cite{lin1988observability},
      \cite{yin2016uniform}, \cite{yin2016synthesis},
      \cite{yin2018synthesis}, 
      \cite{hu2024design}, \cite{xiang2024existence}, 
      \cite{basile2022supervisory}, \cite{aida2020blocking},
      \cite{hu2026supervisory} \\

    & Property-enforcing synthesis
    & \cite{shu2013enforcing}, \cite{ma2022resilient},
      \cite{hu2023asynchronous}, \cite{xie2023optimal},
      \cite{li2025supervisory}, \cite{zheng2026enforcement},
      \cite{zhang2024quantitatively}, \cite{liu2025approximate}, \cite{hu2026supervisory} \\
\bottomrule
\end{tabular}
\end{table}

% 1 partially observable supervisor for the IS-based property
The supervisory control of DESs under partial observation has been extensively studied in the context of information-state-based (IS-based) properties such as safety \cite{ma2022resilient}, diagnosability \cite{li2025supervisory, hu2023asynchronous}, opacity \cite{lafortune2018history, zheng2026enforcement},
detectability \cite{shu2013enforcing}, anonymity and attractability, etc. \cite{yin2024formal, basile2022supervisory}.
These properties characterize whether certain conditions can be inferred from partial observations, and often require synthesis of a supervisor that ensures that all executions of the closed-loop system satisfy the given property.
%
In \cite{yin2016uniform}, a uniform approach is developed to synthesize the bipartite controller for DESs to embed all reachable information states and all admissible supervisory control strategies, which can be further used to synthesize the maximal permissive supervisor for specific properties.
In \cite{xie2023optimal}, this bipartite controller is extended to solve the problem of enforcing opacity for both qualitative (or infinite-step opacity) and quantitative cases.
In \cite{zheng2026enforcement}, the extended version of the bipartite controller is developed for interval weighted DESs for the enforcement of weighted opacity, for which the enabling interval and the sequences of control patterns determine the reachability of the system.
%
However, most of these IS-based properties are defined on prefix-closed specification languages, whereas the non-blocking property requirements are not prefix-closed.
Therefore, simply pruning the illegal states based on the classical bipartite controller is no longer sufficient, which motivates the development of new observation-based progress properties.


% 2 non-blocking problems
Non-blockingness is a fundamental progress requirement in DESs, requiring that every generated behavior can be extended to a marked behavior \cite{chen1991dealing}.
%
The verification of non-blockingness has been extensively investigated for different DES frameworks, including modular DESs \cite{tang2024compositional}, adapted systems \cite{thuijsman2024transformational}, and compositional analysis of extended finite-state machines \cite{mohajerani2016framework}.
%
In \cite{tomas2018complexity}, the computational complexity of verifying nonblockingness in modular DESs is rigorously analyzed, showing that the problem is PSPACE-complete in general, while special cases such as DFA with one shared event or a bounded number of DFAs are NP- or NL-complete.
%
In \cite{pilbrow2015algorithm}, the authors develop a set of compositional abstraction rules, including blocked events, to efficiently verify nonblockingness in modular DESs.
%
In \cite{tang2024compositional}, a framework is developed for compositional non-blocking verification of finite automata with prioritized events, introducing conflict-preserving abstractions and shaped quotient automata to efficiently ensure non-blockingness in modular DESs.
%
However, most of these works on non-blocking analysis focus on verification rather than synthesis. In many practical scenarios, it is often necessary to synthesize a supervisor to actively control the system such that its output behaviors are guaranteed to satisfy the non-blocking property \cite{chen1991dealing, malik2013compositional}. 
This motivates the development of synthesis methods that ensure non-blockingness under various system architectures and observation constraints, bridging the gap between theoretical verification results and practical supervisory control design.


% 3 A + B 
In recent research, the combination of non-blockingness and partial-observation supervisory control has also been widely investigated \cite{fokkink2024offline, hou2023distributed, zhang2024quantitatively}.
%
In \cite{yin2016synthesis}, the IS-based bipartite controller of \cite{yin2016uniform} is extended to synthesize the safe, non-blocking, and maximally permissive supervisor.
%
In \cite{xiang2024existence}, a game-theoretic framework is developed to synthesize nonblocking bounded supervisors for DESs by transforming the supervisor existence problem into a reachability game over information states and admissible control decisions.
%
In \cite{jan2024supervisory}, modular supervisory control under partial observation is investigated from the perspective of normality, where the global nonblocking and maximally permissive normal supervisor is achieved through the parallel composition of local normal supervisors.
%
In \cite{zhang2025bounded},  a new bounded-time nonblocking supervisory control problem is solved for timed DES, in which  each task must be completed within a bounded time counted by the number of ticks.
%
Although significant progress has been made in combining non-blockingness with partial-observation supervisory control \cite{fokkink2024offline, hou2023distributed, zhang2024quantitatively}, non-blockingness only guarantees that the system can eventually reach a marked state (task completion). 
However, it does not ensure that the system can determine from its observations whether it has actually reached the target, leaving a gap in the design of observation-based progress guarantees. 
This motivates the introduction of observation-enforcing non-blocking properties that explicitly connect system observations to task progress.


\subsection{Organization}

The remainder of this article is organized as follows.
%
Section~\ref{sec:prelimi} presents the system model and the notions for partial-observation supervisory control.
Section~\ref{sec:prob} introduces the notion of cognitive non-blockness and gives the problem formulation.
Section~\ref{sec:bcc} illustrates the infeasibility of all-inclusive controller and puts forward the new structure of the  bipartite cognitive controller.
Section~\ref{sec:solu} gives the solution of finding the supervisor that enforces cognitive non-blocking.
Section~\ref{sec:special_z_marked} discuss the special case in which the marked states are within the unobservable reach.
Section~\ref{sec:con} concludes the article and gives the direction for future research.


\section{Preliminaries}\label{sec:prelimi}

This section introduces the preliminaries for the verification and synthesis of cognitive non-blocking in discrete-event systems (DESs), including system models, supervisory control under partial observation, and information states.


\subsection{System model}

We use $\mathbb{N}$, $\mathbb{N}_{\geq 0}$ to denote the sets of natural numbers and nonnegative natural numbers, respectively. 
We use $\varepsilon$ to denote the empty string of an automaton. 
%
The DES discussed is modeled as a \emph{deterministic} finite-state automaton (DFA) and formally defined as follows.
\begin{definition}
    a DFA is a five-tuple
    $$
        G = (Q, \Sigma, \delta, q_0, Q_m)
    $$
    where
    $Q$ is the finite set of states, $\Sigma$ is the finite set of events, 
    $\delta \subseteq Q \times \Sigma \times Q$ is the transition relation,
    $q_0$ is the initial state, and
    $Q_m$ is the set of marked states.
    \label{def:DFA}
\end{definition}
%
In DFA $G$, the \emph{run} is a sequence of states and transitions
$
\rho = q_0 \xrightarrow[]{\sigma_0} q_1 \xrightarrow[]{\sigma_1} \cdots q_n \xrightarrow[]{\sigma_n} q_{n+1},
$
where $(q_{i}, \sigma_i, q_{i+1}) \in \delta$ for $0 \leq i \leq n, n \in \mathbb{N}_{\geq 0}$
The \emph{sequence} corresponding to run $\rho$ is 
$r : tr(\rho) = \sigma_0 \, \sigma_1 \, \cdots \, \sigma_n \in \Sigma^*$,
%
With a slight abuse of notations, we use $\delta(q, r) = q'$ as the partial transitions to denote the reachable state $q'$ from state $q$ and sequence $r$.
%
The behavior generated by $G$ is described by 
$\mathcal{L}(G) = \{ s \in \Sigma^* | \delta(q_0, s)! \}$, 
where $!$ means ``is defined'', and the marked behavior is 
$\mathcal{L}_m(G) = \{ s \in \Sigma^* | \delta(q_0, s) \in Q_m \}$. 
The prefix-closure of a language $\mathcal{L}(G)$ is
$$
\overline{\mathcal{L}(G)} = \{ s \in \Sigma^* | \exists t \in \Sigma^* : st \in \mathcal{L}(G) \},
$$
and we say that $\mathcal{L}$ is prefix-closed if $\overline{\mathcal{L}} = \mathcal{L}$.

In DFA $G$, 
a state $q \in Q$ is said to be \emph{accessible} if there exists a sequence $s \in \mathcal{L}(G)$ such that $\delta(q_0, s) = q$.  
A state $q \in Q$ is \emph{co-accessible} if there exists a sequence $t \in \mathcal{L}(G)$ such that $\delta(q, t) \in Q_m$.

A standard progress property in DES is \emph{non-blocking}, which requires that every generated sequence can be extended to reach a marked state.
\begin{definition}
    Given a DFA $G$, we say that $G$ is \emph{non-blocking} if
    $
    \mathcal{L}(G) = \overline{\mathcal{L}_m(G)}.
    $
    \label{def:nb}
\end{definition}
%
Based on Definition \ref{def:nb}, one can easily check that the system is non-blocking if and only if 
\begin{equation}
    \forall s \in \mathcal{L}(G), \exists t \in \Sigma^* \ s.t. \ st \in \mathcal{L}_m(G).
    \label{eq:nb_1}
\end{equation}
%
Intuitively, this means that from every reachable sequence $s$, the system can still evolve through some continuation $t$ to a marked state; that is, no generated behavior gets trapped before task completion.
Equivalently, in terms of states, non-blocking implies that every state that is \emph{accessible} from the initial state is also \emph{co-accessible}.


We assume the DFA $G$ is partially observed, and the set of events $\Sigma$ can be partitioned into two disjoint subsets, $\Sigma = \Sigma_o \ \dot{\cup} \ \Sigma_{uo}$, where $\Sigma_o$ is the set of observable events, and $\Sigma_{uo}$ is the set of unobservable events.
%
The \emph{natural projection}, $P : \Sigma^* \rightarrow \Sigma_{o}^*$ is recursively defined as
\begin{equation}
P(s \sigma) =
\begin{cases}
P(s), & \text{if } \sigma \in \Sigma_{uo} \\
P(s)\sigma, & \text{otherwise}
\end{cases}
\label{eq:proj}
\end{equation}
with
$P(\varepsilon)  = \varepsilon$.
Intuitively, the projection $P$ prunes the unobservable events within the given sequence $s$ and outputs the sequences of observable events.
For a slight abuse of notations, we denote $P(\mathcal{L}(G))$ as the observable language of DFA $G$.

\subsection{Supervisory control under partial observations}

In the supervisory control framework initiated in \cite{ramadge1987supervisory}, a supervisor is imposed on $G$ to achieve some specification by dynamically enabling or disabling events. 
The event set $\Sigma$ is partitioned into two disjoint subsets: the subset of controllable events $\Sigma_c$, and the subset of uncontrollable events $\Sigma_{uc}$, i.e., 
$\Sigma = \Sigma_c \ \dot{\cup} \ \Sigma_{uc}$.
%
The full set of \emph{control patterns} for DFA $G$ is defined as $\Gamma = 2^{\Sigma_c}$, and the control pattern is $\gamma \in \Gamma$.
%
For DFA $G$, a partial-observation supervisor $S$ is defined as a function $S : P(\mathcal{L}(G)) \rightarrow \Gamma$, such that for $s, t \in \Sigma^*$, we have $S(s) = S(t)$ if $P(s) = P(t)$.
%
We use $S / G$ to denote controlled system, and the language generated by $S /G$, denoted by $\mathcal{L}(S / G)$, is recursively defined recursively as follows,
\begin{enumerate}[(i)]
    \item $\varepsilon \in \mathcal{L}(S / G)$;
    \item $s \, \sigma \in \mathcal{L}(S / G)$ if and only if
    $s \in \mathcal{L}(S / G)$ $\wedge$ $s \, \sigma \in \mathcal{L}(G)$ $\wedge$ $\sigma \in S(s)$.
\end{enumerate}
We also define the marked language of the closed loop system $S / G$, as $\mathcal{L}_m(S / G) = \mathcal{L}(S / G) \cap \mathcal{L}_m(G)$, and
as the \emph{state estimate} upon the occurrence of $s \in \mathcal{L}(G)$, namely, the set of states in which the system could be after observing $P(s) \in P(\mathcal{L}(G))$ as
\begin{equation}\label{eq:esti:1}
\mathcal{E}_S(s) := \{q \in Q | \exists t \in \mathcal{L}(S / G) \ : \ \delta(q_0, t) = q \wedge P(t) = s\}
\end{equation}

\begin{remark}\label{rem:uc_ctrl_pattern}
    In DFAs, in accordance with the Ramadge-Wonham framework of \textup{\cite{ramadge1987supervisory}}, uncontrollable events are always enabled, regardless of the control patterns, that is, transitions labeled by uncontrollable events are guaranteed to be accessible regardless of the control policy.
    %
    To highlight the controllability aspect of the system, we omit uncontrollable events from the control specifications $\gamma$ and $\Gamma$.
    %
    Their accessibility is treated as implicitly enabled in the remainder of this article.
\end{remark}

Let language $K \subseteq \mathcal{L}_m(G)$ be a non-prefix-closed language that represents the desired behavior to be achieved under control. $K$ is said to be:
\begin{enumerate}[(i)]
    \item controllable (w.r.t. DFA $G$ and uncontrollable event set $\Sigma_{uc}$) if $\overline{K} \Sigma_{uc} \cap \mathcal{L}(G) \subseteq \overline{K}$;
    \item observable (w.r.t. DFA $G$, controllable event set $\Sigma_c$ and observable event set $\Sigma_o$) if for all $s \in \overline{K}$ and $\sigma \in \Sigma_c$,
    such that $s \, \sigma \in \overline{K}$, $P^{-1}(P(s)) \, \sigma \cap \mathcal{L}(G) \subseteq \overline{K}$, where $P^{-1}(s) := \{ t \in \Sigma^* | P(t) = s \}$;
    \item $\mathcal{L}_m(G)$-closed if $\overline{K} \cap \mathcal{L}_m(G) = K$. 
\end{enumerate}
%
The following theorem shows that these three conditions together provide the necessary and sufficient conditions for the existence of a supervisor $S$ that exactly achieves the given language $K$.
%
\begin{theorem}[Controllability and Observability Theorem \cite{lin1988observability}]
    Consider DES $G$, with controllable event set $\Sigma_c$ and observable event set $\Sigma_o$. Consider also $K \subseteq \mathcal{L}_m(G)$. 
    There exists a supervisor $S : \mathcal{L}(G) \rightarrow \Gamma$ such that $\mathcal{L}_m(S /G) = K$ and $L(S / G) = \overline{K}$ if and only if
\begin{enumerate}[(i)]
    \item $K$ is controllable w.r.t. $\mathcal{L}(G)$ and $\Sigma_{uc}$;
    \item $K$ is observable w.r.t. $\mathcal{L}(G)$, $\Sigma_o$ and $\Sigma_c$;
    \item $K$ is $\mathcal{L}_m(G)$-closed.
\end{enumerate}
\end{theorem}

If a given specification language cannot be exactly achieved,
then the synthesis problem asks whether we can find a controllable, observable, and $\mathcal{L}_m(G)$-closed sublanguage of the specification language, and one that is ``as large as possible.''
We would like to synthesize the non-blocking and maximally permissive, that is, the closed-loop system $S / G$ satisfies:
\begin{enumerate}[(i)]
    \item $\mathcal{L}_m(S /G) \subseteq K$ and $\mathcal{L}(S /G) = \mathcal{L}_m(S /G)$;
    \item For any other a supervisor $S'$ satisfying (i), we have $\mathcal{L}_m(S / G) \not \subset \mathcal{L}_m(S' / G)$.
\end{enumerate}
%
% Note that (ii) is usually relaxed in the scenario under partial observations, since observability is not preserved under union.


\subsection{Information state and update rules}

\begin{figure*}
    \centering
    \includegraphics[width=0.85\linewidth]{Figs/fsa_aic_overall.pdf}
    \caption{(a) Finite state automaton $G$ with the set of marked states $Q_m = \{ q_3, q_5, q_6, q_{18} \}$, the set of observable events $\Sigma_o = \{ o_1, o_2, o_3 \}$, and the set of controllable events $\Sigma_c = \{ c_1, c_2, o_3 \}$. (b) Bipartite transition system $\mathcal{A}$ constructed from DFA $G$, where the rectangles are for Y-states, and the rounded rectangles are
    for Z-states, respectively, and the control pattern $\emptyset$ in Z-states
    means no events are enabled. The states marked by red crosses are blocking states of Definition V.3 of \cite{yin2016synthesis}, and should be pruned.}
    \label{fig:DFA_aic_overall}
\end{figure*}

Since we are dealing with partially observed systems, we define the notion of an information state as $X = 2^Q$.
We then define the Y-states $Y \subseteq X$ for possible reachable states under partial observations,
and Z-states $Z \subseteq X \times \Gamma$ for states of unobservable reach under supervisory control.
%
The \emph{unobservable reach} of the information states $x \in X$ under
the control pattern $\gamma \subseteq \Sigma$ is given by
\begin{equation}
    \begin{aligned}
    UR_\gamma(x) := \{ &q' \in Q | \exists q \in x, \exists w \in ((\Sigma_{uc} \cup \gamma) \cap \Sigma_{uo})^* \cup \\
    & \{ \varepsilon \} : q' = \delta(q, w) \}.
    \end{aligned}
    \label{eq:ur}
\end{equation}
The \emph{observable reach} of information states $x \in X$ under observable event $\sigma_o \in \Sigma_o$ is given by
\begin{equation}
    NX_{\sigma_o}(x):=\{q'\in Q\mid \exists q\in x : \delta(q,\sigma_o)=q'\}.
    \label{eq:nx}
\end{equation}

To systematically capture the state uncertainty caused by unobservable events and the alternating interaction between control decisions and event occurrences, we adopt an information-state perspective and represent the closed-loop evolution by a bipartite transition system (BTS), which is formally defined as follows.

\begin{definition}
     A bipartite transition system $\mathcal{A}$ with respect to DFA $G$, introduced in \cite{yin2016synthesis}, is a seven-tuple
     $$
     \mathcal{A} = (Y, Z, \delta_{YZ}, \delta_{ZY}, \Sigma, \Gamma, y_0)
     $$
     where
     \begin{enumerate}[(i)]
         \item $Y \subseteq X$ is the set of Y-states;
         %
         \item $Z \subseteq X \times \Gamma$ is the set of Z-states, and
         for simplicity, we use $I(z)$ and $\Gamma(z)$, $z \in Z$ denote the information states and control pattern of a Z-state $z$, such that $z = (x_z, \gamma_z) = (I(z), \Gamma(z))$;
         %
         \item $\delta_{YZ} \subseteq Y \times \Gamma \times Z$ is, for a slight abuse of notation, the partial transition relation from Y-states to Z-states, satisfying
         \begin{equation}
             \delta_{YZ}(y, \gamma) = z \Rightarrow I(z) = UR_{\gamma}(y) \wedge \Gamma(z) = \gamma;
             \label{eq:bts_yz}
         \end{equation}
         %
         \item $\delta_{ZY} \subseteq Z \times \Sigma_o \times Y$ is the partial transition relation from Z-states to Y-states, satisfying
         \begin{equation}
         \begin{aligned}
             \delta_{ZY}(z, \sigma_o) = y \Rightarrow 
             &\sigma_o \in (\Sigma_{uc} \cup \Gamma(z)) \cap \Sigma_o \wedge \\
             &y = NX_{\sigma_o}(I(z));
         \end{aligned}
              \label{eq:bts_zy}
         \end{equation}
        %
        \item $\Sigma$ is the set of events from DFA $G$;
        %
        \item $\Gamma$ is the set of admissible control patterns of DFA $G$;
        %
        \item $y_0$ is the initial Y-state with $y_0 = \{ q_0 \}$.
     \end{enumerate}
    \label{def:bts}
\end{definition}

The notion of BTS is introduced to describe, in a general form, the ``game'' between the supervisor and the system under partial observation. 
%
It is a bipartite structure consisting of two types of states: a Y-state, representing an information state at which the supervisor selects a control decision, and a Z-state, representing the resulting information state together with the selected decision $\Gamma(z)$, from which the system resolves the next observable event.
%
A transition $y \xrightarrow{\gamma} z$ captures the unobservable reach under decision $\gamma$, where $I(z)$ collects the states reachable via enabled unobservable event strings and $\Gamma(z)=\gamma$. 
A transition $z \xrightarrow{\sigma_o} y'$ captures the observable reach after an event $\sigma_o \in \Sigma_o$, where $y'$ consists of the states reachable via $\sigma_o$.
%
A run of the BTS is a sequence $\rho_{\mathcal{A}} = y_0 \xrightarrow{\gamma_0} z_0 \xrightarrow{\sigma_{o,0}} y_1 \cdots$.
We denote the corresponding sequence of states by
$\chi = H(\rho_{\mathcal{A}}) = y_0 z_0 y_1 \cdots$, and denote the sequence of control decisions and observations by 
$s_{\mathcal{A}} = tr(\rho_{\mathcal{A}}) = \gamma_0 \sigma_{o,0} \gamma_1 \sigma_{o,1} \cdots$, 
where $\gamma_i = \Gamma(z_i) \in \Gamma$ and $\sigma_{o,i} \in \Sigma_o$, respectively, for all $i \in \mathbb{N}$.
We use $\bm{q} \in Y \cup Z$ to denote an arbitrary state of the BTS, and $\bm{q}' = \delta_{\mathcal{A}}(\bm{q}, r_{\mathcal{A}})$ to denote the state $\bm{q}'$ transitioned from state $\bm{q}$ via sequence $r_{\mathcal{A}}$.
%
We give the following example to illustrate how the information states evolve under supervisory control.

\begin{example}\label{ex:is_update}
    Let us consider DFA $G$ in Fig. \ref{fig:DFA_aic_overall} (a) with observable event set $\Sigma_o = \{o_1, o_2, o_3\}$ and controllable event set $\Sigma_c = \{ c_1, c_2, o_3 \}$. 
    This DFA is constructed from the motivated example of Fig.~\ref{fig:moti_ex:1}.
    %
    We demonstrate the evolution of the information states using the constructed BTS of \ref{fig:DFA_aic_overall} (b).
    %
    For the initial state $y0 = \{ q_1 \}$, considering $u_1, u_2$ being uncontrollable and unobservable, the only available control pattern is $\gamma_0 = \{ o_3 \}$.
    Next, we have $x_1 = UR_{\gamma_0}(y_0) = \{ q_0, q_1, q_2 \}$ and $z_1 = (x_1, \gamma_1) = (\{ q_0, q_1, q_2 \}, \{ o_3 \})$, where states $q_0$, $q_1$, $q_2$ are accessible from states $q_0$ via unobservable transition $\varepsilon$, $u_1$, and $u_2$, respectively.
    %
    For information state $z_1$, the observable reach based on observation $o_3$ and $o_2$ are
    $y_1 = NX_{o_3}(x_1) = \{ q_3, q_4 \}$ and $y_2 = NX_{o_2}(x_1) = \{ q_5, q_6 \}$.
    Similarly, based on the possible unobservable transition $q_3 \xrightarrow{c_1} q_7 \xrightarrow{c_2} q_9$ and $q_4 \xrightarrow{c_1} q_8 \xrightarrow{c_2} q_{10}$, we can enable nothing $\emptyset$, event $\{ c_1 \}$, event $\{ c_2 \}$ or $\{c_1, c_2 \}$ at this step.
    Then, we have
    $x_4 = UR_{\gamma_2}(y_1) = \{ q_3, q_4, q_7, q_8 \}$ and
    $x_7 = UR_{\gamma_3}(y_1) = \{ q_3, q_4, q_7, q_8, q_9, q_{10} \}$,
    where $\gamma_2 = \{ c_1 \}$ and $\gamma_3 = \{ c_1, c_2 \}$.
    For Z-states $z_7 = (x_7, \gamma_3)$, we have
    $y_5 = NX_{o_2}(x_7) = \{ q_{13}, q_{14} \}$ and
    $y_3 = NX_{o_1}(x_7) = \{ q_{17}, q_{18} \}$.
\end{example}

%
% -----------------------------
\section{Problem formulation}\label{sec:prob}

In this section, we first introduce the notion of cognitive non-blocking based on the preliminaries. 
We then formulate the verification and synthesis problems for DESs under this notion.

\subsection{Non-blocking under partial observations}

In the previous section, Definition~\ref{def:nb} introduces the non-blocking property in DESs.
Under this property, every state that is \emph{accessible} from the initial state $q_0$ is also \emph{co-accessible} to the marked state $Q_m$.
%
The non-blocking supervisory control problem aims to synthesize a supervisor that selects admissible control decisions based solely on partial observations, such that every reachable string in the closed-loop system can be accessible to the marked states. 
In this setting, the supervisor can restrict the behavior of the system by dynamically enabling or disabling the controllable events, thereby rejecting the reachability to the blocking states.

Formally, the non-blocking supervisory control problem under partial observation asks to find a supervisor $S: L(G) \to \Gamma$ such that
\begin{equation}
    \mathcal{L}_m(S / G) \subseteq \mathcal{L}_m(G) \, \wedge \, \mathcal{L}(S / G) = \overline{\mathcal{L}_m(S / G)},
    \label{eq:nb_sc}
\end{equation}
or, in the closed-loop system $S / G$, all states accessible from the initial states are co-accessible to the marked states.
%
This problem is challenging because the supervisor must issue the same control decision for all plant states within an information state, which may include both states that can reach a marked state and states that would lead to blocking. 
Moreover, the number of such information states grows exponentially with the size of the plant, making the synthesis computationally complex.
%
This problem is challenging because partial observation forces the supervisor to make identical control decisions over indistinguishable states, and the maximally permissive supervisor may not exist in this situation.
%

In \cite{yin2016synthesis}, a BTS is introduced as a structural representation of the interaction between the supervisor and the plant under partial observation. 
Define $C_{\mathcal{A}}(y) := \{ \gamma \in \Gamma : \delta_{YZ}(y, \gamma)! \}$ as the set of control decisions defined at $y \in Y$.
Based on the BTS framework, the synthesis problem can be reformulated as the construction of a suitable all-inclusive controller that satisfies the non-blocking requirement, which is defined as follows.
%
\begin{definition}
Given a plant $G$, a BTS $\mathcal{A}$ is called an \emph{all-inclusive controller} (AIC) if it satisfies:
\begin{enumerate}[(i)]
    \item $\mathcal{A}$ is a \emph{complete} BTS, i.e.,
    $
    \forall y \in Y : C_{\mathcal{A}}(y) \neq \emptyset,
    $
    and
    $
    \forall z \in Z, \forall \sigma_o \in \Gamma(z)\cap\Sigma_o:\;
    \big(\exists q \in I(z): \delta(q,\sigma_o)!\big) \Rightarrow \delta_{ZY}(z, \sigma_o)!;
    $
    %
    \item For transition relation $\delta_{YZ}$ and $\delta_{ZY}$, we have
    $$
    \begin{aligned}
        \delta_{YZ}(y, \gamma)   = \, &z \Leftrightarrow (I(z) = UR_{\gamma}(y)) \wedge (\Gamma(z) = \gamma) \\
        \delta_{ZY}(z, \sigma_o) = \, &y \Leftrightarrow (\sigma_o \in \Gamma(z) \cap \Sigma_o) \wedge \\
                                      &(y = NX_{\sigma_o}(I(z))).
    \end{aligned}
    $$
\end{enumerate}
\label{def:aic}
\end{definition}
%
Note that we omit the explicit construction of the safety specification here, as it is beyond the scope of this article.
In condition (i), for any reachable Y-state, there should be at least one control decision at that state, and for all enabled observable events at a Z-state, their corresponding transitions should be defined if they exist in DFA $G$.
In condition (ii), the BTS only satisfies the ``if'' part, which is different from \eqref{eq:bts_yz} and \eqref{eq:bts_zy}.

In AIC $\mathcal{A}$, each evolution corresponds to a sublanguage of $\mathcal{L}(G)$.
Let $\mathcal{L}(\mathcal{A})$ denote the language generated by the AIC $\mathcal{A}$, and denote $L \in \mathcal{L}(\mathcal{A})$ as a language generated by the AIC $\mathcal{A}$, 
i.e., the set of all sequences in $\Sigma^*$ that can be realized under some supervisor included in $\mathcal{A}$.
%
The following theorem provides us with the solution to the problem of synthesizing non-blocking supervisory controllers based on partial observation.

\begin{theorem}[Theorem V.1 \cite{yin2016synthesis}]\label{theo:nb_aic}
Given an AIC $\mathcal{A}$ constructed from DFA $G$, $\mathcal{A}$ is said to be \emph{non-blocking} if it satisfies the following conditions,
\begin{enumerate}[(i)]
    \item the AIC $\mathcal{A}$ is \emph{complete};
    \item every reachable $Y$-state is \emph{live}, i.e., for every reachable
    $y\in Y$ and every $q\in y$, there exist $n \in \mathbb{N}$ and a decision
    string $\gamma_1 \gamma_2 \cdots \gamma_n \in \Gamma^n$ with
    $\gamma_1\in C_{\mathcal{A}}(y)$ such that the following condition holds:
    there exists a string
    $$
        s=\xi_1 \sigma_1 \xi_2 \sigma_2 \cdots \xi_{n-1} \sigma_{n-1} \xi_n \in \Sigma^*
    $$
    for DFA $G$
    satisfying $\xi_i \in ((\gamma_i \cup \Sigma_{uc}) \cap \Sigma_{uo})^*$
    and $\sigma_i \in (\gamma_i \cap \Sigma_o) \cup (\Sigma_{uc} \cap \Sigma_o)$,
    $\forall i \in \mathbb{N},$
    such that
    $
        \delta(q,s) \in Q_m;
    $
    % Moreover, by letting $y_0=y$ and
    % $
    %     y_i=\delta_{ZY}(\delta_{YZ}(y_{i-1},\gamma_i),\sigma_i),
    % $
    % we require that each $y_i$ is well-defined and
    % $
    %     \gamma_{i+1}\in C_{\mathcal{A}}(y_i), \forall i \in \mathbb{N};
    % $
    %
    \item every reachable $Z$-state is \emph{deadlock-free}, i.e., for every
    reachable $z\in Z$ and every $q\in I(z)$, we have    
    $$
    \begin{aligned}
    &(\exists s \in (\Gamma_a(z) \cap \Sigma_{uo})^* \Rightarrow \delta(q, s) \in Q_m) \vee \\
    &(\exists s \in (\Gamma_a(z) \cap \Sigma_{uo})^*(\Gamma_a(z) \cap \Sigma_o) \Rightarrow \delta(q, s)!),
    \end{aligned}
    $$
    where $\Gamma_a(z) = \Gamma(z) \cup \Sigma_{uc}$.   
\end{enumerate}
Then, the language generated by the non-blocking AIC (NB-AIC) $\mathcal{L}(\mathcal{A})$ satisfies
\begin{enumerate}[(i)]
    \item If $L \in \mathcal{L}(\mathcal{A})$, then $L$ is controllable, observable and deadlock-free.
    \item If $L$ is non-empty, controllable, observable, and non-blocking, then $L \in \mathcal{L}(\mathcal{A})$.
\end{enumerate}
\end{theorem}
Intuitively, the liveness of $Y$-states requires that, for each plant state
contained in the current information state, there exists a sequence of
admissible control patterns under which this plant state can reach a marked
state. 
The deadlock-free property requires that the Z-states be accessible to the successor live Y-state in the sense that the marked states can be eventually reached.
%
Now, we give the following example to illustrate how to synthesize an NB-AIC from a given DFA.

% subsection 1
% 1 non-blocking is widely considered
% 2 BTS can be used to solve the NB problem (definition)
% 3 however, NB is insufficient

\begin{example}
    Let us consider the DFA $G$ and the corresponding BTS $\mathcal{A}$ illustrated in Fig.~\ref{fig:DFA_aic_overall} again.
    One can easily check that this BTS is an AIC for the Definition \ref{def:aic}.
    Then we can construct the NB-AIC by pruning those blocking information states according to Theorem \ref{theo:nb_aic}.
    %
    First, for Z-states 
    $z_2 = (\{ q_{3},  q_{4}  \}, \emptyset)$,
    $z_3 = (\{ q_{5},  q_{6}  \}, \emptyset)$,
    $z_6 = (\{ q_{13}, q_{14} \}, \emptyset)$, and
    $z_{12} = (\{ q_{15}, q_{16} \}, \emptyset)$ are deadlock Z-states, so they should be pruned.
    %
    Next, based on the completeness condition of Y-states, at least one control decision should be given at Y-states, so we prune $y_5 = \{ q_{13}, q_{14} \}$ for the pruning $z_6$.
    Then, considering that all possible observations should be illustrated for any Z-states, we prune
    $z_7 = (\{ q_3, q_4, q_7, q_8,    q_9,    q_{10} \}, \{ c_1, c_2 \})$ and
    $z_8 = (\{ q_5, q_6, q_9, q_{10}, q_{11}, q_{12} \}, \{ c_1, c_2 \})$, such that the remaining AIC is non-blocking.
    %
    Let us take state $q_{4}$ for example, based on the evolution of AIC $\mathcal{A}$
    \begin{equation}\label{eq:ex_nb:1}
        \{ o_3 \} \, o_3 \, (\{ c_1 \} \, o_1 \, \emptyset \, o_1)^\omega,
    \end{equation}
    and the marked states $q_{3}$ and $q_{18}$ is reachable via the run 
    \begin{flalign}
    \rho_1 &= q_0 \xrightarrow{u_1} q_1  \xrightarrow{o_3} (q_3 \xrightarrow{c_1} q_7 \xrightarrow{o_1} q_{17} \xrightarrow{o_1} q_{3} )^\omega \label{eq:ex_nb:2} \\
    %
    \rho_2 &= q_0 \xrightarrow{u_2} q_2  \xrightarrow{o_3} (q_4 \xrightarrow{c_1} q_8 \xrightarrow{o_1} q_{18} \xrightarrow{o_1} q_{4} )^\omega,  \label{eq:ex_nb:3}
    \end{flalign}
    where $(\cdot)^\omega$ means the sequence will be repeated infinitely many times.
    %
    Similar checks can be done for the rest of the states of DFA $G$, verifying that the controlled system is non-blocking.
    \label{ex:nb_1}
\end{example}


\subsection{Problem formulation for cognitive non-blocking}

From \eqref{eq:nb_sc} and \eqref{eq:ex_nb:1}, in the closed-loop system $S/G$, for each state $q \in Q$ reached by $s$, there exists a finite sequence of admissible events $\sigma_1 \cdots \sigma_k \in \Sigma^*$ 
leading from $q$ to some marked state $q_m \in Q_m$, i.e.,
\begin{equation}
    \forall s \in \mathcal{L}(S/G), \exists t \in \Sigma^* \ s.t. \ st \in \mathcal{L}_m(S/G).
    \label{eq:nb_sc:2}
\end{equation}
%
However, the non-blockingness only illustrates the capability of the states within DFA $G$ for arriving at a marked state.
In this sense, if the system only has partial observability of the occurrence of events and does not know the exact state at each step, then it may not be able to determine for sure the arrival of the marked state.
%
We give the following example to illustrate that, under partial observation, the supervisor may not be able to determine whether the system is non-blocking based solely on its observations.
%
\begin{example}\label{ex:cnb_1}
    Consider the DFA $G$ from Example~\ref{ex:nb_1}, with marked states $Q_m = \{ q_3, q_5, q_6, q_{18} \}$. 
    When applying the control sequence \eqref{eq:ex_nb:1}, the system run can follow either \eqref{eq:ex_nb:2} or \eqref{eq:ex_nb:3}. 
    From the AIC in Fig.~\ref{fig:bcc_overall}(b), the observable reachable sets are $y_1 = \{ q_3, q_4 \}$ or $y_3 = \{ q_{17}, q_{18} \}$. 
    In other words, based solely on observations, the observer cannot distinguish which of the two runs the system has taken.
    For example, once observing $y_1$ cannot determine whether the current state is either $q_3$ or $q_4$.
    Consequently, although both runs satisfy classical non-blockingness, the observer cannot definitively determine that the system has reached a marked state.
    % 这个写的不好，就是无论如何观测都
\end{example}
%
Therefore, we would like to discuss the notion named \emph{cognitive non-blocking}, in which the partial observation can determine for sure the non-blockingness based on its partial observations, 
\begin{equation}
    \forall s \in \mathcal{L}(S/G), \exists t \in \Sigma^* \Rightarrow \mathcal{E}_S(st) \subseteq Q_m,
\end{equation}
%
where $\mathcal{E}_S(st)$ is defined in \eqref{eq:esti:1}.
Intuitively, a DFA satisfies cognitive non-blocking if it can eventually certify, based on its observations, that it is in a marked state. 
That is, there exists a continuation such that the resulting information state consists entirely of marked states.
% It is straightforward to verify that cognitive non-blocking is a stronger property than classical non-blocking.

Before giving the problem formulation, we would like to discuss the physical interpretation behind this definition.
Cognitive non-blocking emphasizes the ability to \emph{certify} task completion based on observations.
%
However, if a marked state is reachable from a nonmarked state via unobservable events, the system may enter the marked set silently, making it impossible for an external observer to determine whether the task has indeed been completed.
Therefore, interpreting reachability of marked states as task completion is meaningful only when such completion can be inferred from observations.
% 
In many physical scenarios, such as a robot arriving at a goal region or a machine finishing an operation, completion is expected to be accompanied by some observable event.
%
To this end, we refine the notion of cognitive non-blocking and formally formulate the problem as follows.
%
\begin{definition}[Cognitive Non-blocking]
    Given DFA $G = (Q, \Sigma, \delta, q_0, Q_m)$ and a set of observable events $\Sigma_o \subseteq \Sigma$, we say DFA $G$ is \textit{cognitive non-blocking} with respect to observable event set $\Sigma_o$ if
    \begin{equation}\label{eq:cnb_1}
    \forall s \in \mathcal{L}(S/G), \exists t \in \Sigma^*, \Rightarrow \hat{\mathcal{E}}_S(st) \subseteq Q_m,
    \end{equation}
    where
    \begin{equation}\label{eq:est_o}
    \begin{aligned}
    \hat{\mathcal{E}}_S(s) := \{ 
    &q \in Q : \exists t \in (\Sigma^*\Sigma_o) \cap \mathcal{L}(S / G) \, s.t. \\ 
    &\delta(q_0, t) = q \wedge P(t) = P(s) \}
    \end{aligned}
    \end{equation} 
    is the \emph{cognitive state estimate} associated with sequence $s \in P(\mathcal{L}(G))$, i.e., the set of states that can be reached immediately after the immediate observation of $s$.       
    \label{def:cnb}
\end{definition}
%
\begin{problem}[Cognitive non-blocking synthesis problem]
    Given DFA $G = (Q, \Sigma, \delta, q_0, Q_m)$, a set of observable events $\Sigma_o \subseteq \Sigma$, and a set of controllable events $\Sigma_c \subseteq \Sigma$, synthesize the partial-observation supervisor such that the system is \textit{cognitive non-blocking} with respect to the marked states accessible via the observable events
    \begin{equation}
    \forall s \in \mathcal{L}(S/G), \exists t \in \Sigma^*, \Rightarrow \hat{\mathcal{E}}_S(st) \subseteq Q_m.
    \end{equation}
    \label{prob:cnb}
\end{problem}






% subsection 2
% 1 what is CNB (definition)
% 2 why BTS cannot solve CNB? (proposition)
% 3 give the problem


\section{Bipartite cognitive controller under all-inclusive controller}\label{sec:bcc}

In this section, we first show the structure of AIC is not capable of solving Problem~\ref{prob:cnb}. 
To this end, we introduce the bipartite cognitive controller (BCC) that captures the microscopic evolution within the BTS such that the cognitive non-blocking property can be verified.

\subsection{Insufficiency of all-inclusive controller}

We begin by presenting the following lemma, illustrating that cognitive non-blocking is stronger than classical non-blocking.

\begin{lemma}
    Given DFA $G$, a set of observable events $\Sigma_o \subseteq \Sigma$, and a set of controllable events $\Sigma_c \subseteq \Sigma$, if $G$ is said to be cognitive non-blocking, it is non-blocking, but not vice versa.
    \label{lem:cnb_nb}
\end{lemma}
\begin{proof}
    Based on \eqref{eq:nb_1} and \eqref{eq:est_o}, we show that cognitive non-blocking implies non-blocking, or
    $\hat{\mathcal{E}}_S(st) \subseteq Q_m$ is a subset of $\mathcal{L}_m(G)$.
    With 
    $$
    \begin{aligned}
    \hat{\mathcal{E}}_S(st) = \{& q \in Q : \exists t' \in (\Sigma^*\Sigma_o) \cap \mathcal{L}(S/G), \\
                                &\delta(q_0,t') = q,\ P(t') = st \} \, \subseteq \, Q_m
    \end{aligned}
    $$ 
    implying that for all $t'$ satisfying $P(t') = st$, we have
    $
    \delta(q_0, t') \in Q_m.
    $
    In particular, there exists at least one such $t' \in \mathcal{L}(S/G)$, which implies
    $
    t' \in \mathcal{L}_m(G).
    $
    Therefore, for any $t' \in \mathcal{L}(S/G)$, there exists a continuation leading to a marked state, and thus $G$ is non-blocking.
    %
    The converse does not hold. That is, from \eqref{eq:nb_sc:2}, there exists a DFA $G$ such that
    $$
        \forall s \in \mathcal{L}(S/G), \exists t \in \Sigma^* \ s.t. \delta(q_0, st) \in Q_m,
    $$
    but cognitive non-blocking means
    $$
        \forall s \in \mathcal{L}(S/G), \ \exists t \in \Sigma^*, \ 
        \hat{\mathcal{E}}_S(st) \subseteq Q_m,
    $$
    where $\hat{\mathcal{E}}_S(\cdot)$ requires that $st \in (\Sigma^*\Sigma_o)$ should be observation-consistent sequences rather than universal sequences.
\end{proof}

Based on Lemma~\ref{lem:cnb_nb}, it is straightforward to see that Example~\ref{ex:cnb_1} illustrates a case where the AIC satisfies the classical non-blocking property, but fails to satisfy cognitive non-blocking.

We now present the following proposition to show that Theorem~\ref{theo:nb_aic} is insufficient to guarantee cognitive non-blockingness.

\begin{proposition}\label{prop:insuff_aic}
Given an AIC $\mathcal{A}$ constructed from DFA $G$, if it satisfies 
    (i) $\mathcal{A}$ is complete, 
    (ii) every reachable $Y$-state is live, and 
    (iii) every reachable $Z$-state is deadlock-free,
then $\mathcal{A}$ is non-blocking, but it does not guarantee cognitive non-blocking.
\end{proposition}
%
\begin{proof}
The non-blockingness of the AIC $\mathcal{A}$ is guaranteed by Theorem~\ref{theo:nb_aic}, 
so we only need to show that the cognitive non-blocking property may fail.

Assume, for contradiction, that $\mathcal{A}$ satisfies both non-blocking and cognitive non-blocking. 
Based on the liveness requirement of Y-states in Theorem~\ref{theo:nb_aic} and the observability of events leading to marked states, 
for every $q \in y$, there exists a string
$$
    s = \xi_1 \sigma_1 \xi_2 \sigma_2 \cdots \xi_{n-1} \sigma_{n-1} \xi_n \sigma_n \in \Sigma^*\Sigma_o
$$
such that $\delta(q,s) \in Q_m$, where
$\xi_i \in ((\gamma_i \cup \Sigma_{uc}) \cap \Sigma_{uo})^*$ 
and $\sigma_i \in ((\gamma_i \cup \Sigma_{uc}) \cap \Sigma_o)$ for all $i$.
%
However, there can exist a system execution where a sequence $t'$ satisfies $P(t') = P(s)$ 
but $\delta(q,t') \notin Q_m$. In such a case, the cognitive state estimate satisfies
$
    \hat{\mathcal{E}}_S(s) \not\subseteq Q_m,
$
which contradicts the assumption of cognitive non-blocking.
%
Therefore, although $\mathcal{A}$ is non-blocking in the classical sense, it may exhibit cognitive blocking, 
and the proposition holds.
\end{proof}

From the proof of Proposition~\ref{prop:insuff_aic}, we observe that the fundamental reason for the DFA's infeasibility is that the liveness of Y-states is insufficient to meet the stronger cognitive non-blocking requirement. 
While the deadlock-free property of Z-states ensures classical non-blocking within unobservable reaches, the liveness of Y-states does not guarantee that the completion of the tasks can be confirmed based on the cognitive state estimate.

% 
% 然后需要说明为什么不足够解决，参考一下博涵的聊天记录
% 原文中的NB是直接用AIC做的
% Z状态是deadlock-free是for all states的(公式(7))，而Y-状态是live的
% 这里似乎是不是Z的条件是一致的，似乎我们只要加强Y状态即可？
% 那么我的问题变成了是不是只用AIC就可以解决这个问题 (就是动机这关怎么过)
%   答案是:
%   Z是不变的，但是就算判断live也要用ICS (BCC)
%   如果没有 ICS/BCC 结构，你就无法直接判断 Y-state 是否 live，因为普通 NB-AIC 的 live 判断是针对真实状态，而 cognitive live 需要追踪信息状态 evolution（即观察到的 projection 下的可能状态集合）。
%
% 对于某个 Y-state y（对应 NB-AIC 中的信息状态集合），它是 cognitive live，当且仅当对于每个 x∈y 的可能真实状态，都存在一条控制决策序列，使得经过一系列允许的事件后，observer 能确认系统进入了 fully marked Y-state。
%
% 3. ICS/BCC 结构的必要性
% ICS/BCC（Information Consistency Structure / Bipartite Control Consistency）结构本质上是：
% 展开信息状态：对每个 Y-state，追踪在 observable projection 下，可能的 plant states 集合 evolution。
% 可控决策序列：记录哪些控制决策保证从该信息状态出发可以到达 fully marked Y-state。
% 递归判定 live：通过 BFS/DFS 或 fixed-point，判断每个信息状态是否有至少一条控制序列使 cognitive estimate 最终落入 fully marked Y-state。
% 因此：
% 没有 ICS/BCC，就无法考虑信息状态 evolution，只能用 plant state 判定 live；
% cognitive live 必须依赖 ICS/BCC 或等价结构，否则 supervisor 无法保证从 observer 角度看非阻塞。
% 换句话说：
%       cognitive live \neq = plant live (原 NB-AIC live)
% 只有通过信息状态展开（ICS/BCC），才能把 live 条件提升到 observer 可验证的层面。

%
% 这个模板caption中\bm{q}显示不出来
\begin{figure*}[t!]
    \centering
    \includegraphics[width=0.675\linewidth]{Figs/bcc_overall.pdf}
    \caption{Bipartite cognitive controller constructed from the DFA $G$ and the AIC $\mathcal{A}$ in Fig.~\ref{fig:DFA_aic_overall}. We simplify the notations $(y,q)$ and $(z,q)$ to $q$ to better illustrate the structural properties of each microstate. Rectangles represent Y-states, while rounded rectangles represent Z-states. Transitions in different colors are used to distinguish different evolutions. Moreover, the information states marked by red crosses are infeasible Z-states, cognitive-blocking Y-states, or inconsistent Y- or Z-states, which should be removed to achieve cognitive non-blocking. 
    The source code is available at \url{https://github.com/trigger1996/des_common}.}

    \label{fig:bcc_overall}
\end{figure*}

\subsection{Bipartite cognitive controller under all-inclusive controller}

To address this issue that relying solely on the NB-AIC framework fails to address cognitive live conditions, we need a structure that tracks the microscopic evolution of information states under observable events, rather than just based on observations. 
The bipartite cognitive controller (BCC) structure fulfills this role by expanding information states and recording admissible control sequences.
Now, we give the formal definition of BCC as follows.
%
\begin{definition}\label{def:bcc}
Given an AIC $\mathcal{A}$, its corresponding
BCC is a tuple
$
\mathcal{C} = (\bm{\mathcal{Q}}, \Sigma_\mathcal{C}, \delta_\mathcal{C}, q^Y_0, \bm{\mathcal{Q}}_m)
$
where
\begin{enumerate}[(i)]
    \item $\bm{\mathcal{Q}} \subseteq \bm{\mathcal{Q}}^Y \dot{\cup} \bm{\mathcal{Q}}^Z$ is the set of microstates defined by
    $$
    \begin{aligned}
        \bm{\mathcal{Q}}^Y &= \{ (y, q) \in Y \times Q \mid y \in Y \wedge q \in y \} \\
        \bm{\mathcal{Q}}^Z &= \{ (z, q) \in Z \times Q \mid z \in Z \wedge q \in I(z) \},        
    \end{aligned}
    $$

    \item $\Sigma_\mathcal{C} = \Sigma \cup \Gamma$ is the set of label of transitions.

    \item $\delta_\mathcal{C} \subseteq \bm{\mathcal{Q}} \times \Sigma_\mathcal{C} \times \bm{\mathcal{Q}}$ is the transition relation:
    \begin{itemize}
        \item $\delta_\mathcal{C}((y,q_1),\gamma) = (z,q_2)$ if $q_1=q_2$, $\delta_{YZ}(y,\gamma)=z$, $\gamma=\Gamma(z)$,
        \item $\delta_\mathcal{C}((z,q_1),\sigma) = (z,q_2)$ if $\delta(q_1,\sigma)=q_2$ and $\sigma \in (\Gamma(z)\cup\Sigma_{uc}) \cap \Sigma_{uo}$,
        \item $\delta_\mathcal{C}((z,q_1),\sigma) = (y,q_2)$ if $\delta(q_1,\sigma)=q_2$, $\sigma \in (\Gamma(z)\cup\Sigma_{uc}) \cap \Sigma_o$, and $\delta_{ZY}(z,\sigma)=y$.
    \end{itemize}

    \item $q^Y_0 = (\{q_0\}, q_0)$ is the initial state.

    \item $\bm{\mathcal{Q}}_m = \{ (y,q) \in \bm{\mathcal{Q}}^Y \mid y\subseteq Q_m, q\in y \}$ is the set of cognitively marked microstates.
\end{enumerate}
\end{definition}
%
In BCC $\mathcal{C}$, we use $q^Y = (y, q) \in \bm{\mathcal{Q}}_y$ and $q^Z = (z, q) \in \bm{\mathcal{Q}}_z$ to denote the individual microstates.
Given Y-state $y \in Y$ and Z-state $Z \in Z$ of an AIC $\mathcal{A}$, its corresponding microstates, $\bm{\mathcal{Q}}_y$ and $\bm{\mathcal{Q}}_z$ is defined as 
$
\bm{\mathcal Q}_y := \{ (y,q) \in \bm{\mathcal Q}^Y \mid q \in y \}
$
and
$
\bm{\mathcal Q}_z := \{ (z,q) \in \bm{\mathcal Q}^Z \mid q \in I(z) \},
$
respectively.
%
A \emph{run} of $\mathcal{C}$ is a sequence of microstates
$
\rho_\mathcal{C} = q_0^Y q_{1,1}^Z \cdots q_{1,m_1}^Z q_1^Y q_{2,1}^Z \cdots,
$
and we denote the corresponding \emph{sequence} of events and control patterns for BCC $\mathcal{C}$ as
$$
\bm{s} = tr_\mathcal{C}(\rho_\mathcal{C}) = \gamma_0 \sigma_{uo,1,1} \cdots \sigma_{uo,1,m_1} \sigma_{o,1} \gamma_1 \cdots \in \Sigma_\mathcal{C}^*,
$$
where we use the boldface symbol such as $\bm{s}$ to explicitly record these sequences, distinguishing them from ordinary event sequences of DFAs.
%
Recall that $\bm{q}$ is used to denote $Y \cup Z$.
Following the definition of the transition function $\delta$ for DFA $G$, we define the partial transition function $(\bm{q}', q') = \delta_\mathcal{C}((\bm{q}, q), tr_\mathcal{C}(\rho_\mathcal{C}))$, such that $(\bm{q}, q)$ and $(\bm{q}', q')$ are the origining and ending microstates of the run $\rho_\mathcal{C})$, respectively.
We also define the full set of the runs of BCC as $\mathcal{R}_\mathcal{C}$.
%
Intuitively, each run of $\mathcal{C}$ is a microstate-level refinement of the corresponding evolution of $\mathcal{A}$, and the first component of state $q^Y$ or $q^Z$ gives the Y/Z-state trajectory, and the second component gives the underlying trajectory in $G$.
%
Following this idea, we establish the relation from the runs of AIC $\mathcal{A}$ to the runs of BCC $\mathcal{G}$.
Given the run of BCC $\rho_\mathcal{C} = (y_0, q_0) (z_1, q_{1, 1}) \cdots (z_1, q_{1, m}) (y_1, q_{y, 1}) (z_2, q_{2, 1}) \cdots$, its corresponding run of AIC is $\rho_{\mathcal{A}} = y_0 \xrightarrow{\gamma_0} z_1 \xrightarrow{\sigma_{o, 0}} y_1 \cdots$.
%
% 下到上
% \highlight{HERE WE NEED TO REFINE THE STATE SEQUENCES RATHER THAN RUNS, JUST SEE $\{ q_{15}, q_{16} \}$ can reach $\{q_0\}$ via $o_3$ and $o_1$.}
We use operator $R_{\mathcal{C}}^\mathcal{A} : \mathcal{R}_{\mathcal{C}} \rightarrow X^*$ to obtain the information state sequence from the run $\rho_{\mathcal{C}}$, such that
\begin{flalign}
    R_{\mathcal{C}}^\mathcal{A}(\rho_\mathcal{C}) = \{ 
    &y_0 z_1 y_1 z_2 \cdots \in X^* \mid \notag\\
    & \rho_\mathcal{C} = q_0^Y q_{1,1}^Z \cdots q_{1,m_1}^Z q_1^Y q_{2,1}^Z \cdots, \ \notag\\
    &q_i^Y = (y_i, q), \ q_{i,j}^Z = (z_i, q') \ \forall i,j \}. \notag
\end{flalign}
%
Moreover, we also define operator $R_{\mathcal{A}}^\mathcal{C} : X^* \rightarrow 2^{\mathcal{R}_{\mathcal{C}}}$  for information state sequence $\chi = H(\rho_\mathcal{A})= y_0 z_1 y_1 z_2 \cdots$, such that
$$
\begin{aligned}
    R_{\mathcal{A}}^\mathcal{C}(\chi) := \{ &\rho_{\mathcal{C}}' \in \mathcal{R}_\mathcal{C} \mid 
        \rho_{\mathcal{C}}' = (y_0, q_0) (z_1, q_{1,1}') \cdots (z_1, q_{1,m}') (y_1, q_{y,1}') \cdots, \\
        & q_{1,1}',\ldots,q_{1,m}' \in I(z_1), \ q_{y,1}' \in y_1, \ldots \}.
\end{aligned}
$$
Note that the AIC run $\rho_{\mathcal{A}}$ may correspond to multiple BCC runs $R_{\mathcal{C}}^\mathcal{A}(\rho_{\mathcal{C}})$.
%
The language generated by BCC $\mathcal{C}$ is defined by
\begin{equation}
    \mathcal{L}(\mathcal{C}) := \{ tr_{\mathcal{C}}(\rho_{\mathcal{C}}) \in \Sigma_{\mathcal{C}}^* \mid \delta_{\mathcal{C}}(q^Y_0, tr_{\mathcal{C}}(\rho_{\mathcal{C}}))! \}.
    \label{eq:lang_bcc}
\end{equation}
Since $\mathcal{C}$ contains both plant events and control patterns, we define the projection of events
$
    P_\Sigma: \Sigma_{\mathcal{C}}^* \to \Sigma^*
$
for BCC $\mathcal{C}$ recursively by
\begin{equation}\label{eq:seq_proj_bcc}
P_\Sigma(\alpha \eta)
=
\begin{cases}
    P_\Sigma(\alpha) \eta, & \eta \in \Sigma,\\
    P_\Sigma(\alpha),      & \eta \in \Gamma,
\end{cases}
\end{equation}
with $P_\Sigma(\varepsilon)=\varepsilon$,
$\forall \alpha \in \Sigma_{\mathcal{C}}^*$ and $\eta \in \Sigma_{\mathcal{C}}$.
Following the projection $P_{\Sigma}$ in \eqref{eq:seq_proj_bcc}, the projection onto control patterns is defined recursively as
\begin{equation}\label{eq:proj_bcc_ctrl}
P_\Gamma(\alpha \eta) =
\begin{cases}
P_\Gamma(\alpha) \eta, & \eta \in \Gamma,\\
P_\Gamma(\alpha), & \eta \in \Sigma,
\end{cases}
\quad
P_\Gamma(\varepsilon)=\varepsilon.
\end{equation}
%
Finally, we define the projected plant language induced by $\mathcal{C}$ as
\begin{equation}
    \mathcal{L}_{\Sigma}(\mathcal{C})
    :=
    P_\Sigma(\mathcal{L}(\mathcal{C}))
    \subseteq \Sigma^* .
    \label{eq:proj_lang_bcc}
\end{equation}
Since $\mathcal{L}_{\Sigma}(\mathcal{C})$ is generated by an automaton-based structure, it is naturally prefix-closed.

\begin{algorithm}[ht!]
\caption{Construction of BCC}
\label{alg:bcc}
\KwIn{DFA $G = (Q, \Sigma, \delta, q_0, Q_m)$\;
      $\qquad$ $\;$
      AIC $\mathcal{A} = (Y, Z, \delta_{YZ}, \delta_{ZY}, \Sigma, \Gamma, y_0)$\;}
\KwOut{BCC $\mathcal{C} = (\bm{\mathcal{Q}}, \Sigma_\mathcal{C}, \delta_\mathcal{C}, q^Y_0, \bm{\mathcal{Q}}_m)$\;}

Initialize $\bm{\mathcal{Q}} \gets \{(y_0, q_0)\}$, $\delta_\mathcal{C} \gets \emptyset$, $Y\_stack \gets \{ (y_0,q_0)\}$, $visited \gets \{(y_0,q_0)\}$\;

\While{$Y\_stack \neq \emptyset$}{
    Pop $(y,q)$ from $Y\_stack$\;

    \ForEach{$(y, \gamma, z) \in \delta_{YZ}$}{
        % \tcp{Y->Z control transition}
        \If{$(z,q) \notin visited$}{
            Add $(z,q)$ to $\bm{\mathcal{Q}}$\;
            Add edge $((y,q), \gamma, (z,q))$ to $\delta_\mathcal{C}$\;
            $visited \gets visited \cup \{(z,q)\}$\;
        }

        % \tcp{Internal unobservable transitions in Z, propagate step by step}
        $Z\_stack \gets [(z,q)]$\;
        \While{$Z\_stack \neq \emptyset$}{
            Pop $(z,q_{curr})$ from $Z\_stack$\;
            \ForEach{$\sigma \in (\Gamma(z) \cup \Sigma_{uc}) \cap \Sigma_{uo}$}{
                $q' \gets \delta(q_{curr}, \sigma)$\;
                \If{$q' \in I(z) \wedge (z,q') \notin visited$}{
                    Add $(z,q')$ to $\bm{\mathcal{Q}}$\;
                    Add transition $((z,q_{curr}), \sigma, (z,q'))$ to $\delta_\mathcal{C}$\;
                    $visited \gets visited \cup \{(z,q')\}$\;
                    Push $(z,q')$ to $Z\_stack$\;
                }
            }
        }

        % \tcp{Observable exits Z->Y}
        \ForEach{$(z, \sigma_o, y') \in \delta_{ZY}$}{
            $q' \gets \delta(q,\sigma_o)$\;
            \If{$(y',q') \notin visited$}{
                Add $(y',q')$ to $\bm{\mathcal{Q}}$\;
                Add transition $((z,q), \sigma_o, (y',q'))$ to $\delta_\mathcal{C}$\;
                Push $(y',q')$ to $Y\_stack$\;
                $visited \gets visited \cup \{(y',q')\}$\;
            }
        }
    }
}

$\bm{\mathcal{Q}}_m \gets \{ (y,q) \in \bm{\mathcal{Q}} \mid y \in Y \wedge q \in Q_m \}$\;
\Return $(\bm{\mathcal{Q}}, \Sigma_\mathcal{C}, \delta_\mathcal{C}, (y_0,q_0), \bm{\mathcal{Q}}_m)$\;
\end{algorithm}

We give Algorithm~\ref{alg:bcc} for the construction of a BCC from the original DFA $G$ and AIC $\mathcal{A}$.
% 
Algorithm~\ref{alg:bcc} follows a breadth-first search (BFS) structure.
It iteratively finds all the successor Z-microstates from the given Y-microstates and finds the successor Y-microstates of the new Z-microstates until all states are visited.
%
In Lines 2 -- 9, we perform a BFS over Y-microstates: for Y-microstate $(y,q)$, we explore all Y-Z control transitions in $\delta_{YZ}$.
For each new Z-microstate $(z,q)$, we add it to the BCC and create the Y-Z edge representing the control decision.
%
In Lines 10--18, we propagate Z-microstates internally through unobservable events: each $(z,q_{curr})$ generates new Z-microstates $(z,q')$ for each unobservable event $\sigma$ and recursively pushes them to the Z-stack.
%
In Lines 19--26, we handle observable exits from Z to Y.
For each observable event $\sigma_o$, we compute $q' = \delta(q, \sigma_o)$ and add new Y-microstate $(y',q')$ and the corresponding edge labeled by $\sigma_o$ to the BCC.
We use Example~\ref{ex:bcc} to illustrate the construction of BCC.
\begin{example}\label{ex:bcc}
    Consider again the DFA $G$ and its AIC $\mathcal{A}$ in Fig.~\ref{fig:DFA_aic_overall}. 
    We construct the corresponding BCC using Algorithm~\ref{alg:bcc}, with the resulting structure shown in Fig.~\ref{fig:bcc_overall}, where rectangles and rounded rectangles denote Y-states and Z-states, respectively. 
    For clarity, we simplify the notation $(\bm{q}, q)$ to $q$ and use different colors for the transition relations to highlight distinct evolutions, which better illustrate the structural properties of each microstate.
\end{example}

\subsection{Properties of bipartite cognitive controller}

The BCC represents a microstate-level expansion of an AIC. 
It captures both the information states induced by the partially observable supervisor and the detailed transitions among individual DFA states, including those occurring within unobservable reaches and those induced by observable events from a given state.


Consider a BCC $\mathcal{C} = (\bm{\mathcal{Q}}, \Sigma_\mathcal{C}, \delta_\mathcal{C}, q^Y_0, \bm{\mathcal{Q}}_m)$, where 
each state is of the form $(\bm{q}, q)$ with $\bm{q} \in Y \cup Z$ and $q \in Q$, and each transition is of the form $((\bm{q}, q), \sigma, (\bm{q}', q'))$. 
Define
\begin{flalign}\label{eq:pre_post}
    Pre((\bm{q}, q))  = \{ &(\bm{q}', q') \in \bm{\mathcal{Q}} \mid \exists \sigma \in \Sigma_{\mathcal{C}}, \nonumber \\ 
    &((\bm{q}', q'), \sigma, (\bm{q}, q)) \in \delta_{\mathcal{C}} \} \nonumber \\
    %
    Post((\bm{q}, q)) = \{ &(\bm{q}', q') \in \bm{\mathcal{Q}} \mid \exists \sigma \in \Sigma_{\mathcal{C}}, \nonumber \\ 
    &((\bm{q}, q), \sigma, (\bm{q}', q')) \in \delta_{\mathcal{C}} \}
\end{flalign}
as the sets of predecessors and successors of $(\bm{q}, q)$, respectively. 
For a set of vertices $\bm{\mathcal{Q}}_{s} \subseteq \bm{\mathcal{Q}}$, we use
\begin{flalign}\label{eq:pre_post:2}
    Pre(\bm{\mathcal{Q}}_{s})  := &\bigcup_{(\bm{q}, q) \in \bm{\mathcal{Q}}_{s}} Pre((\bm{q}, q)) \nonumber \\
    Post(\bm{\mathcal{Q}}_{s}) := &\bigcup_{(\bm{q}, q) \in \bm{\mathcal{Q}}_{s}} Post((\bm{q}, q))
\end{flalign}
to denote the sets of predecessors and successors, respectively. 
%
Let $\bm{\mathcal{Q}}_{t} \subseteq \bm{\mathcal{Q}}$ denote the set of target vertices. 
For $(\bm{q}, q), (\bm{q}', q') \in \bm{\mathcal{Q}}$ and $\bm{\mathcal{Q}}_s \subseteq \bm{\mathcal{Q}}$, we write $(\bm{q}', q') \leadsto_{\bm{\mathcal{Q}}_s} (\bm{q}, q)$ if there exists a nonempty path from $(\bm{q}', q')$ to $(\bm{q}, q)$,
$$
(\bm{q}', q') (\bm{q}_{1}, q_1) (\bm{q}_{2}, q_2) \cdots (\bm{q}_{n}, q_n) (\bm{q}, q),
$$
such that $(\bm{q}_{1}, q_1), \ldots, (\bm{q}_{n}, q_n) \in \bm{\mathcal{Q}}_s$.
A state $(\bm{q}, q) \in \bm{\mathcal{Q}}$ is called a \emph{deadlock} if 
\begin{equation}\label{eq:bcc_dl}
(\bm{q}, q) \notin \bm{\mathcal{Q}}_{t} \wedge Post((\bm{q}, q)) = \emptyset,
\end{equation}
a nonempty set $\bm{\mathcal{Q}}_l \subseteq \bm{\mathcal{Q}}$ is called a \emph{livelock} if $\bm{\mathcal{Q}}_l \cap \bm{\mathcal{Q}}_t = \emptyset$, $\bm{\mathcal{Q}}_l$ is a nontrivial strongly connected component, and $Post(\bm{\mathcal{Q}}_l) \subseteq \bm{\mathcal{Q}}_l$, i.e.,
\begin{flalign}\label{eq:bcc_ll}
&\bm{\mathcal{Q}}_l \cap \bm{\mathcal{Q}}_t = \emptyset \wedge
\forall (\bm{q},q), (\bm{q}',q') \in \bm{\mathcal{Q}}_l : 
((\bm{q},q) \leadsto_{\bm{\mathcal{Q}}_l} \nonumber \\
& (\bm{q}',q') \wedge
  (\bm{q}',q') \leadsto_{\bm{\mathcal{Q}}_l} (\bm{q},q)) \wedge
Post(\bm{\mathcal{Q}}_l) \subseteq \bm{\mathcal{Q}}_l.   
\end{flalign}
%
These notions can be naturally extended to other automata models, including the DFA $G$ and AIC $\mathcal{A}$, by interpreting $\bm{\mathcal{Q}}$ and $\delta_\mathcal{C}$ as the corresponding sets of states and transitions, respectively.

Based on the above notations, we analyze the properties of the BCC, illustrating the consistency of its evolution with the original DFA, the completeness of the control representation, and its construction complexity in both time and space.
%
We first discuss the structure of Z-states in an AIC $\mathcal{A}$ and their corresponding Z-microstates $\{ (z, q) \mid q \in I(z) \}$.
Given a Z-microstate $(z, q)$, if all its predecessor states $Pre((z,q))$ belong to $\bm{\mathcal{Q}}^Y$, then it is called an \emph{entry Z-microstate}.
If $Post((z,q)) \cap \bm{\mathcal{Q}}^Y \neq \emptyset$, it is called an \emph{exit Z-microstate}.
Note that a Z-microstate can be both an entry and an exit microstate.
%
Given an entry Z-microstate $(z, q)$ and any of its predecessor Y-microstates $(y, q)$, considering that in \eqref{eq:ur}, $\varepsilon$ can always be used as the unobservable transition such that $q = \delta(q, \varepsilon)$, it holds that $q \in I(z) \cap y$.
The transitions among Z-microstates within the same Z-state correspond to unobservable reaches from their predecessor Y-states.
Given an entry Z-microstate $(z, q)$, for every $q' \in I(z)$, there exists a string $s \in \Sigma_{uo}^*$ such that $(z, q') = \delta_{\mathcal{C}}((z, q), s)$.
%
For Y-states, they represent observation-equivalent states under supervisory control.
Thus, for any $y \in Y$ and any $(y, q), (y, q')$ with $q, q' \in y$, we have $((y, q), \xi, (y, q')) \notin \delta_{\mathcal{C}}$ for all $\xi \in \Sigma_{\mathcal{C}}$.
%
Then, we give the following proposition to illustrate the consistency between the DFA $G$ and the BCC $\mathcal{C}$.
%
\begin{proposition}\label{prop:bcc_dfa_consistency}
Given a run 
$$
\begin{aligned}
\rho_{\mathcal{C}} = &(y_0, q_0) (z_1, q_{1,1}) \cdots (z_1, q_{1,m_1}) (y_1, q_1) \cdots \\
                     &(y_i, q_i) (z_{i+1}, q_{i+1,1}) \cdots (z_{i+1}, q_{i+1,m_{i+1}}) (y_{i+1}, q_{i+1}) \cdots
\end{aligned}
$$
in the BCC $\mathcal{C}$, there exists a corresponding sequence
$
s = \xi_0 \sigma_{o,0} \xi_1 \sigma_{o,1} \cdots \xi_i \sigma_{o,i} \cdots
$
in the closed-loop system $S/G$ such that:
\begin{enumerate}[(i)]
    \item Each $\xi_i \in ((\Gamma(z_i) \cup \Sigma_{uc}) \cap \Sigma_{uo})^*$ is a unobservable subsequence corresponding to the unobservable reach $(z_{i}, q_{i,1}) \cdots (z_{i}, q_{i,m_{i}}) (y_{i}, q_{i})$;
    \item For each subsequence $(z_i, q_{i,1}) \cdots (z_i, q_{i,m_i})$, there exist events $\sigma_{i,j} \in \Sigma_{uo}$ such that
    $
    (q_{i,j}, \sigma_{i,j}, q_{i,j+1}) \in \delta;
    $
    \item For each transition $(z_i, q_{i,m_i}), \sigma_{o,i}, (y_{i+1}, q_{i+1})$, we have
    $
    (q_{i,m_i}, \sigma_{o,i}, q_{i+1}) \in \delta.
    $
\end{enumerate}
\end{proposition}
%
\begin{proof}
The result follows directly from the construction of $\delta_{\mathcal{C}}$.
%
First, every Z-Z transition $((z, q), \sigma, (z, q')) \in \delta_{\mathcal{C}}$ is defined only if $\sigma \in (\Gamma(z) \cup \Sigma_{uc}) \cap \Sigma_{uo}$ and $(q, \sigma, q') \in \delta$. 
Thus, each subsequence $(z_i, q_{i,1}) \cdots (z_i, q_{i,m_i})$ corresponds to a sequence $\xi_i \in ((\Gamma(z_i) \cup \Sigma_{uc}) \cap \Sigma_{uo})^*$ satisfying (ii).
%
Second, every Z-Y transition $((z, q), \sigma_o, (y, q')) \in \delta_{\mathcal{C}}$ is defined only if $\sigma_o \in \Sigma_o$ and $(q, \sigma_o, q') \in \delta$, which establishes (iii).
%
Finally, every Y-Z transition $((y_i, q_i), \gamma_{i+1}, (z_{i+1}, q_{i+1,1}))$ is defined by $\delta_\mathcal{C}$ such that $\gamma_{i+1} = \Gamma(z_{i+1})$.
This determines the set of admissible events $(\Gamma(z_{i+1}) \cup \Sigma_{uc})$ that govern the subsequent unobservable evolution within $z_{i+1}$, and thus constrains the corresponding sequence $\xi_{i+1}$.
Concatenating all such unobservable sequences $\xi_i$ and observable events $\sigma_{o,i}$ yields a valid execution $s$ in the closed-loop system $S/G$, which proves (i).
\end{proof}
%
As a direct consequence of Proposition~\ref{prop:bcc_dfa_consistency}, 
for any $\alpha \in \mathcal{L}(\mathcal{C})$, if
$
\delta_{\mathcal{C}}(q^Y_0,\alpha)=(\bm{q},q),
$
then
$
    \delta(q_0,P_\Sigma(\alpha)) = q.
$
In particular, we have
$$
    \mathcal{L}_{\Sigma}(\mathcal{C})
    =
    P_\Sigma(\mathcal{L}(\mathcal{C}))
    \subseteq
    \mathcal{L}(G).
$$
Based on the above propositions, we now show that the behavior induced by the BCC is closed under all plant-feasible uncontrollable events, and hence is controllable.
%
% Controllability保证 pruning 后的 BCC 不是一个“理想化行为集合”，而是一个在 uncontrollable events 下封闭、可以由 supervisor 实际实现的 closed-loop language。
\begin{theorem}\label{theo:ctrl_bcc}
Given BCC 
$\mathcal{C} = (\bm{\mathcal{Q}}, \Sigma_{\mathcal{C}}, \delta_{\mathcal{C}}, q^Y_0, \bm{\mathcal{Q}}_m)$
constructed from a deterministic DFA
$G=(Q,\Sigma,\delta,q_0,Q_m)$
and AIC
$\mathcal{A}=(Y,Z,\delta_{YZ},\delta_{ZY},\Sigma,\Gamma,y_0)$.
%
Then the projected plant language induced by $\mathcal{C}$,
$
    \mathcal{L}_{\Sigma}(\mathcal{C}) = P_\Sigma(\mathcal{L}(\mathcal{C})),
$
is controllable with respect to $G$ and $\Sigma_{uc}$, i.e.,
\begin{equation}
    \mathcal{L}_{\Sigma}(\mathcal{C})\Sigma_{uc} \cap \mathcal{L}(G)\subseteq \mathcal{L}_{\Sigma}(\mathcal{C}).
    \label{eq:bcc_controllability}
\end{equation}
\end{theorem}
%
\begin{proof}
It suffices to prove that for any 
$s \in \mathcal{L}_{\Sigma}(\mathcal{C})$ and any 
$\sigma \in \Sigma_{uc}$, we have
$$
    s\sigma \in \mathcal{L}(G) \Rightarrow s\sigma \in \mathcal{L}_{\Sigma}(\mathcal{C}).
$$
Let $s \in \mathcal{L}_{\Sigma}(\mathcal{C})$ and
$\sigma \in \Sigma_{uc}$ be such that
$
    s\sigma \in \mathcal{L}(G).
$
By the definition of $\mathcal{L}_{\Sigma}(\mathcal{C})$, there exists
$\alpha \in \mathcal{L}(\mathcal{C})$ such that
$
    P_\Sigma(\alpha) = s.
$
Let
$
    \delta_{\mathcal{C}}(q^Y_0,\alpha)=(\bm{q},q).
$
By the construction of the BCC, every transition labeled by an event follows
the transition relation of DFA $G$, while every control pattern in $\Gamma$
does not change the plant component. Hence,
$$
    q = \delta(q_0,P_\Sigma(\alpha)) = \delta(q_0,s).
$$
Since $s\sigma\in\mathcal{L}(G)$, there exists $q'\in Q$ such that
$
    \delta(q,\sigma) = q'.
$

To prove $s\sigma \in \mathcal{L}_{\Sigma}(\mathcal{C})$, it suffices to
construct a string $\beta \in \Sigma_{\mathcal{C}}^*$ such that
$
    \delta_{\mathcal{C}}((\bm{q},q),\beta)!
$
and
$
    P_\Sigma(\alpha\beta) = s\sigma,
$
Then, we may have $\alpha \beta \in \mathcal{L}(\mathcal{C})$, and
$s\sigma \in \mathcal{L}_{\Sigma}(\mathcal{C})$.
%
First, suppose that $\bm{q} = y \in Y$. By the completeness of AIC
$\mathcal{A}$, there exists an available control decision
$\gamma \in C_{\mathcal{A}}(y)$ and $z\in Z$ such that
$
    \delta_{YZ}(y,\gamma) = z.
$
By the definition of $\delta_{\mathcal{C}}$, we have
$
    \delta_{\mathcal{C}}((y,q),\gamma) = (z,q).
$
Since $\gamma \in \Gamma$, it is pruned by $P_\Sigma$, with
$$
    P_\Sigma(\alpha\gamma) = P_\Sigma(\alpha) = s,
$$
the run reaches the Z-microstate $(z,q)$.
%
If $\beta = \gamma \sigma$ and $\sigma \in \Sigma_{uo}$, then we have
$$
    \sigma\in\Sigma_{uc}\cap\Sigma_{uo} \subseteq (\Gamma(z)\cup\Sigma_{uc})\cap\Sigma_{uo}.
$$
Together with $\delta(q,\sigma) = q'$, the BCC transition rule gives
$
    \delta_{\mathcal{C}}((z,q),\sigma) = (z,q'),
$
such that
$
    \delta_{\mathcal{C}}((y,q),\beta)!
$
and
$$
    P_\Sigma(\alpha\beta) = P_\Sigma(\alpha\gamma\sigma) = P_\Sigma(\alpha)\sigma = s \sigma.
$$
%
If $\beta = \gamma \sigma$ and $\sigma \in \Sigma_o$, since $\sigma \in \Sigma_{uc}$ and $(z,q) \in \bm{\mathcal{Q}}$, we have
$$
    \sigma\in\Sigma_{uc}\cap\Sigma_o \subseteq (\Gamma(z)\cup\Sigma_{uc})\cap\Sigma_o,
$$
and $q\in I(z)$. 
Together with $\delta(q,\sigma) = q'$, the completeness of AIC $\mathcal{A}$ implies that there exists $y'\in Y$ such that
$
    \delta_{ZY}(z,\sigma) = y'.
$
By the definition of $\delta_{\mathcal{C}}$, we have
$
    \delta_{\mathcal{C}}((z,q),\sigma) = (y',q').
$
Thus, by taking $\beta = \gamma\sigma$, we have
$
    \delta_{\mathcal{C}}((y,q), \beta)!
$
and
$$
    P_\Sigma(\alpha\beta) = P_\Sigma(\alpha\gamma\sigma) = P_\Sigma(\alpha)\sigma = s\sigma.
$$

Second, suppose that $\bm{q} = z\in Z$. Then the current microstate is already
a Z-microstate.
%
If $\sigma \in \Sigma_{uo}$, then we have
$$
    \sigma\in\Sigma_{uc}\cap\Sigma_{uo} \subseteq
    (\Gamma(z)\cup\Sigma_{uc})\cap\Sigma_{uo}.
$$
Together with $\delta(q,\sigma) = q'$, the BCC transition rule gives
$
    \delta_{\mathcal{C}}((z,q),\sigma) = (z,q').
$
Thus, by taking $\beta = \sigma$, we have
$
    \delta_{\mathcal{C}}((z,q),\beta)!
$
and
$$
    P_\Sigma(\alpha\beta) = P_\Sigma(\alpha\sigma) = P_\Sigma(\alpha)\sigma = s\sigma.
$$
%
If $\sigma\in\Sigma_o$, then we have 
$$
    \sigma\in\Sigma_{uc}\cap\Sigma_o \subseteq (\Gamma(z)\cup\Sigma_{uc})\cap\Sigma_o.
$$
Since $(z,q)\in\bm{\mathcal{Q}}$, we have $q\in I(z)$. Together with
$\delta(q,\sigma) = q'$, the completeness of AIC $\mathcal{A}$ implies that
there exists $y'\in Y$ such that
$
    \delta_{ZY}(z,\sigma) = y'.
$
By the definition of $\delta_{\mathcal{C}}$, we have
$
    \delta_{\mathcal{C}}((z,q),\sigma) = (y',q').
$
Thus, by taking $\beta = \sigma$, we have
$
    \delta_{\mathcal{C}}((z,q),\beta)!
$
and
$$
    P_\Sigma(\alpha\beta) = P_\Sigma(\alpha\sigma) = P_\Sigma(\alpha)\sigma = s\sigma.
$$

Since $\Sigma = \Sigma_o \, \dot{\cup} \, \Sigma_{uo}$, the above cases cover all
possible uncontrollable events $\sigma\in\Sigma_{uc}$. Therefore,
$s\sigma\in\mathcal{L}_{\Sigma}(\mathcal{C})$. Hence,
$$
    \mathcal{L}_{\Sigma}(\mathcal{C})\Sigma_{uc} \cap \mathcal{L}(G) \subseteq \mathcal{L}_{\Sigma}(\mathcal{C}),
$$
which proves that $\mathcal{L}_{\Sigma}(\mathcal{C})$ is controllable with
respect to $G$ and $\Sigma_{uc}$.
\end{proof}

The significance of Theorem~\ref{theo:ctrl_bcc} is that the projected plant language of the pruned BCC is closed under all plant-feasible uncontrollable events. 
Therefore, it can be implemented as a closed-loop language by a supervisor, since the plant cannot be forced to leave $\mathcal{L}_{\Sigma}(\mathcal{C})$ through the occurrence of an uncontrollable event.
%
Moreover, we provide the time complexity of Algorithm~\ref{alg:bcc} and the space complexity of BCC as follows.

\begin{proposition}
Given a DFA $G$ with $|Q|$ states and $|\Sigma|$ events, Algorithm~\ref{alg:bcc} runs in time
$
    O(2^{|Q|+|\Sigma|}\cdot |Q|^2\cdot |\Sigma|).
$
\end{proposition}

\begin{proof}
Algorithm~\ref{alg:bcc} searches over the microstates of the BCC.
%
In Line 2, the outer loop explores Y-microstates $(y,q)$. Since each
$y\in Y$ contains at most $|Q|$ plant states, the number of possible
Y-microstates is at most $|Y|\cdot |Q|$.

In Line 4, for each Y-microstate $(y,q)$, the algorithm traverses outgoing
Y-Z transitions $(y, \gamma, z)\in\delta_{YZ}$. Since each control decision
$\gamma$ is a subset of $\Sigma_c$, there are at most $2^{|\Sigma_c|}$
possible decisions.
%
In Lines 10--12, for each selected transition $(y, \gamma, z)$, the algorithm
propagates within the corresponding Z-state. For a fixed $z$, there are at
most $|Q|$ Z-microstates. For each Z-microstate, the algorithm checks events
in
$
    (\Gamma(z) \cup \Sigma_{uc}) \cap \Sigma_{uo},
$
whose cardinality is at most $|\Sigma|$. Hence, the number of repetitions is
$
    O(|Q|\cdot|\Sigma|)
$
for each Y-Z transition.
%
In Line 19, the algorithm scans observable exits
$(z, \sigma_o, y') \in \delta_{ZY}$. For each $z$, there are at most
$|\Sigma_o| \le |\Sigma|$ such exits, which is dominated by
$
    O(|Q| \cdot |\Sigma|)
$.
%
Combining the above loops, the total running time is bounded by
$$
    O\big(
    |Y| \cdot |Q| \cdot 2^{|\Sigma_c|}
    \cdot (2 \cdot |Q| \cdot |\Sigma|)
    \big)
    =
    O(|Y|\cdot |Q|^2\cdot |\Sigma|\cdot 2^{|\Sigma_c|}).
$$
Since $|Y| \le 2^{|Q|}$ and $|\Sigma_c| \le |\Sigma|$, we obtain
$$
    O(2^{|Q|+|\Sigma|} \cdot |Q|^2\cdot |\Sigma|).
$$
\end{proof}

%
\begin{proposition}\label{prop:state_bcc}
Consider a DFA $G$ with $|Q|$ states and $|\Sigma|$ events.
For the corresponding BCC $\mathcal{C}$, the number of microstates satisfies
$$
    |\bm{\mathcal{Q}}| \leq |Q|\big(2^{|Q|}+2^{|Q|+|\Sigma|}\big),
$$
and the number of transitions satisfies
$$
    |\delta_{\mathcal{C}}| \leq |Q|2^{|Q| + \Sigma}(1 + |\Sigma|).
$$
\end{proposition}

\begin{proof}
Each BCC microstate is of the form $(\bm{q},q)$, where
$\bm{q}\in Y\cup Z$ and $q\in Q$. Hence,
$
    |\bm{\mathcal{Q}}| \le (|Y|+|Z|)\cdot |Q|.
$
Since $Y\subseteq 2^Q$ and $Z\subseteq 2^Q\times 2^{\Sigma_c}$, we have
$
|Y|\le 2^{|Q|}
$
and
$
|Z|\le 2^{|Q|+|\Sigma_c|}
\le
2^{|Q|+|\Sigma|}.
$
Therefore,
$$
    |\bm{\mathcal{Q}}|
    \le
    |Q|\big(2^{|Q|}+2^{|Q|+|\Sigma|}\big).
$$

The transitions of the BCC are of three types at the microstate level:
Y-Z transitions, Z-Z transitions, and Z-Y transitions.
Since the underlying plant is deterministic, each microstate has at most
one successor under each fixed event or control decision.
Thus, the numbers of Y-Z, Z-Z, and Z-Y micro-transitions can be bounded by
the numbers of corresponding microstates multiplied by the numbers of
available control decisions, unobservable events, and observable events,
respectively.
Thus,
$$
    |\delta_{\mathcal{C}}|
    \le
    |Y||Q| \cdot 2^{|\Sigma_c|}
    +
    |Z||Q| \cdot |\Sigma_{uo}|
    +
    |Z||Q| \cdot |\Sigma_o|.
$$
Using
$
|Y|\le 2^{|Q|}
$,
$
|Z|\le 2^{|Q|+|\Sigma_c|}
$,
and
$
|\Sigma_o|+|\Sigma_{uo}|=|\Sigma|,
$
we obtain
$$
\begin{aligned}
|\delta_{\mathcal{C}}| 
&\leq |Q|2^{|Q| + |\Sigma_c|} + |Q||\Sigma|2^{|Q| + |\Sigma_c|} \\
&\leq |Q|2^{|Q| + \Sigma}(1 + |\Sigma|)
\end{aligned}
$$
\end{proof}



\section{Solution to Cognitive non-blocking synthesis problem}\label{sec:solu}
% 
% 我后面需要什么，我前面需要什么
% 
%
% 1 consistency of BCC 作为下一章节的第一个定理
% 2 what is potentially feasible Z-state?
% 3 what is cognitive live Y-state
% 4 how to synthesis CNB
% 5 Algorithm 2
% 6 correctness of Algorithm 2 (有必要用单源单汇路径的结论吗？)
% 7 

In this section, we leverage the definition of consistency for BCCs.
By pruning those inconsistent states in BCC, the sublanguage generated by the BCC can be guaranteed to be cognitive non-blocking.
%

\subsection{Synthesis algorithm}\label{subsec:supervisor_induction}

To synthesize cognitive non-blocking sublanguages, we first identify the
target microstates characterized by \eqref{eq:cnb_1}.

Recall that in Definition~\ref{def:bcc}, the marked states $\bm{\mathcal{Q}}_m$ are defined as the set of Y-microstates such that for each $q^Y \in \bm{\mathcal{Q}}_m$, its corresponding Y-state $y_m \in Y$ consists entirely of marked plant states. 
For each such cognitively marked Y-state $y_m$, we denote its corresponding Y-microstate set as $\bm{\mathcal{Q}}_{y,m}$.
%
Intuitively, according to \eqref{eq:cnb_1}, a target Y-state corresponds to an information state in which all observation-consistent plant states are marked. 
Therefore, regardless of which actual plant state has been reached under the partial observation, it can be certified that the system has reached a marked state.
%
Meanwhile, we define the consistency of a BCC $\mathcal{C}$ as follows.
%
\begin{definition}\label{def:con_bcc}
Given a BCC $\mathcal{C}$, we say it is \emph{consistent} if and only if its corresponding AIC $\mathcal{A}$ is consistent, that is,
\begin{enumerate}[(i)]
    \item For each Y-microstate $(y, q)$, at least one control decision is defined for its information state $y \in Y$, 
    $
    \forall (y, q) \in \bm{\mathcal{Q}}^Y, y \in Y : C_{\mathcal{A}}(y) \neq \emptyset
    $. 
    \item For each Z-microstate $(z, q)$, all possible observable transitions from its information state $z \in Z$ is defined, 
    $
    \forall (z, q) \in \bm{\mathcal{Q}}^Z, z \in Z, \forall \sigma_o \in \Gamma(z)\cap\Sigma_o:\;
    \big(\exists q \in I(z): \delta(q,\sigma_o)!\big) \Rightarrow \delta_{ZY}(z, \sigma_o)!;    
    $
    % TODO to check

\end{enumerate}
\end{definition}
%
One can easily check that the BCC is consistent if and only if its corresponding AIC is consistent.
%
Next, we give the following two definitions for properties of Y-states and Z-states based on the reachability of their microstates, which will be further utilized for the notion of cognitive non-blocking.
%
\begin{definition}
    We say a Z-state $z \in Z$ is \emph{potentially feasible} if, for each of its microstates $q^Z \in \bm{\mathcal{Q}}_z$, it does not lead to a deadlock \eqref{eq:bcc_dl} or a livelock \eqref{eq:bcc_ll} consisting solely of the Z-microstates $\bm{\mathcal{Q}}_z$ of $z$. 
\end{definition}
%
% 这里有一点很关键，这个Y-state必须能到同一个mark，不能(y, q)到一个marked, (y, q')到另一个mark
\begin{definition}
    We say a Y-state $y$ is \emph{cognitively live} with respect to the set of marked Y-microstates $\bm{\mathcal{Q}}_m$ if and only if there exists cognitively marked Y-state $y_m \subseteq Q_{m}$ such that all Y-microstates $(y,q)$ satisfy
    \begin{flalign}\label{eq:bcc_y_cl}
        &\forall q \in y, \, \exists \chi, \exists q_m \in y_m : \rho_{\mathcal{C}} \in R_{\mathcal{A}}^{\mathcal{C}}(\chi) \Rightarrow \notag \\
        &(y_m, q_m) = \delta_{\mathcal{C}}((y,q), tr(\rho_{\mathcal{C}})),
    \end{flalign}
    where $\chi$ is an arbitrary state sequence of AIC $\mathcal{A}$ via potentially feasible Z-microstates and other cognitively live Y-states.
\end{definition}
%
Intuitively, a Y-state is cognitively live if all its microstates can reach the microstates of cognitively marked Y-state via the runs $R_{\mathcal{A}}^{\mathcal{C}}(\chi)$ corresponding to the state sequence of AIC, $\chi$, of the potentially feasible Z-microstates and other cognitively live Y-states.
%
% CRITICAL
If a Z-state is potentially feasible, then all of its microstates are potentially feasible, and vice versa.
Moreover, if a Y-state is said to be cognitively live, then its Y-microstates are cognitively live.
%
We give the following example to illustrate in detail the property of potential feasibility and cognitive liveness in a BCC.


\begin{figure}[ht!]
    \centering
    \includegraphics[width=\linewidth]{Figs/bcc_routes_overall.pdf}
    \caption{Illustration of several evolutions from the BCC $\mathcal{C}$ in Fig.~\ref{fig:bcc_overall}. For simplicity, the Z-states corresponding to plant state $q_{17}$ can be either $(\{ q_{15}, q_{16}, q_{17} \}, \{ c_1 \})$ or $(\{ q_{15}, q_{16}, q_{17} \}, \{ c_1, o_3 \})$. That is, the enabling of event $o_3$ determines whether the Y-microstate $(\{q_0\}, q_0)$ is accessible from the Z-microstate of $q_{17}$, and this transition is highlighted in orange.}
    \label{fig:bcc_routes}
\end{figure}
%
\begin{example}\label{ex:y_z_prop}
    Let us consider again the DFA $G$ in Fig.~\ref{fig:DFA_aic_overall} (i) and the corresponding BCC $\mathcal{C}$ in Fig.~\ref{fig:bcc_overall}.
    %
    Two of its evolutions are illustrated in Fig.~\ref{fig:bcc_routes}.
    %
    In the evolution of Fig.~\ref{fig:bcc_routes} (a), the Z-state $z_6 = (\{ q_{13}, q_{14} \}, \emptyset)$ is not potentially feasible, for the all its Z-microstates $(z_6, q_{13})$ and $(z_6, q_{14})$ are deadlock states.
    %
    Meanwhile, the Y-state $y_5 = \{ q_{13}, q_{14} \}$ is not cognitively live, for its only available control decision is to deadlock Z-states $z_6$.
    As for Y-states $y_1 = \{ q_{3}, q_{4} \}$ and $y_3 = \{ q_{17}, q_{18} \}$, they are not cognitively live, either.
    With the only available marked Y-state being $y_2 = \{ q_5, q_6 \}$, there is no available sequence from Y-micostates  $(y_1, q_3)$, $(y_1, q_4)$, $(y_3, q_{17})$, $(y_4, q_{18})$ to marked states to $(y_2, q_5)$ and $(y_2, q_6)$, so \eqref{eq:bcc_y_cl} will not be satisfied.
    %
    As for the evolution of Fig.~\ref{fig:bcc_routes} (b), Y-state $y_4 = \{ q_{15}, q_{16} \}$ is cognitively live.
    For state sequence $\chi = y_4 z_{11} y_0 z_0$, with $z_{11} = (\{ q_{15}, q_{16}, q_{17}\}, \{c_1, o_3\})$, $y_0 = \{q_0\}$, and $z_0 = (\{ q_0, q_1, q_2 \}, \emptyset)$, we have the corresponding BCC runs
    $
    \rho_{\mathcal{C}, 1} = (y_4, q_{15}) \xrightarrow{\{c_1, o_3\}} (z_{11}, q_{15}) \xrightarrow{c_1} (z_{11}, q_{17}) \xrightarrow{o_3} (y_0, q_0) \xrightarrow{\emptyset} (z_0, q_0) \xrightarrow{u_1} (z_0, q_1) \xrightarrow{o_2} (y_2, q_5)
    $
    $
    \rho_{\mathcal{C}, 2} = (y_4, q_{16}) \xrightarrow{\{c_1, o_3\}} (z_{11}, q_{16}) \xrightarrow{o_1} (y_0, q_0) \xrightarrow{\emptyset} (z_0, q_0) \xrightarrow{u_2} (z_0, q_2) \xrightarrow{o_2} (y_2, q_6)
    $
    ending at microstates of identical cognitive marked states $y_2$.
    %
    Similarly, one can easily check that $(y_0, q_0)$ is also cognitively live.
    Moreover, Z-states $z_{11} = (\{ q_{15}, q_{16}, q_{17}\}, \{c_1, o_3\})$ and $z_5 = (\{q_5, q_6, q_{11}, q_{12}\}, \{ c_1 \})$ are potentially feasible, for the successor of exit Z-states, $(z_5, q_{11})$, $(z_5, q_{12})$, $(z_{11}, q_{16})$, and $(z_{11}, q_{17})$ are cognitively live Y-microstates.
    However, if the controllable event $o_3$ is not enabled at $z_2$, then the run $\rho_{\mathcal{C}, 1}$ is no longer feasible.
    In this sense, $y_4$ is no longer cognitive live.
    %  
    % That is, by observing the subsequence $(o_3 o_2)^\omega$ or $(o_1 o_2)^\omega$, the observer can certify that the system arrives at a marked state.
    %
    % 这个问题在于15、16到q_0的trace不一样，一个o_1，一个o_3, 所以不能是run到run，必须是state sequence到run
    %
\end{example}

%
% 1   考虑Y和Z是耦合的，而Z直接影响系统的可行性，Y只需要保证有通路即可
% 2   首先标记那些不可行的Z-state，通过Z-microstates判断
% 2.2 此时infeasible Z可以分为两类：一类是all所有结束状态为deadlock
% 2.3                             一类是存在exit states
% 3   首先删除2.2和2.3
% 4   再根据consistency倒着删，此时一边确定系统是否consistent，一边确定Y是否cognitively live，直到系统CNB
%
%
% 我知道最后的结果就是系统会很破碎，但是：不行的状态就是不行
To synthesize a supervisor that guarantees the sublanguage of the closed-loop system is cognitive non-blocking, we need to prune those states that do not satisfy the properties of potentially feasible, cognitively live, or cognitively marked.
Moreover, based on Definition~\ref{def:con_bcc}, when pruning these microstates, we have to prune the corresponding microstates that share identical Y or Z states.
%
Although the accessibility of Y-states and Z-states is coupled, which determines the non-blocking property of the system during its evolution, 
any Z-microstate that is a deadlock or livelock within an unobservable reach is clearly infeasible. 
Thus, we start by identifying Z-states that are not potentially feasible. 
Specifically, we first detect the Z-microstates corresponding to deadlocks and livelocks, as defined in \eqref{eq:bcc_dl} and \eqref{eq:bcc_ll}, and then prune the associated Z-states from the BCC.
%
The removal of infeasible Z-states may affect the consistency and cognitive liveness of predecessor Y-states. 
For each predecessor Y-state $y \in Y$ of a pruned Z-state, 
if $y$ is no longer cognitively live -i.e., \eqref{eq:bcc_y_cl} does not hold for some $q \in y$, or it violates Definition~\ref{def:con_bcc} (i) — then $y$ should also be pruned.
%
After updating Y-states, some remaining Z-states may become inconsistent. 
We recursively prune inconsistent Z-states, and their predecessor cognitively live Y-states until no further removals are necessary and the system becomes consistent. 

We summarize this procedure in Algorithm~\ref{alg:cnb_syn} and illustrate it with Example~\ref{ex:cnb_removal}.
%
% TODO
% 这部分的正确性的concern的担忧主要来自于删除顺序，因为这个问题比原有的IS-based问题在删除过程中会有更强的耦合性，所以得看怎么删
% 最好的办法是过一遍代码检查一下
Note that the resulting Y-states after the pruning procedure are not guaranteed to be accessible from the initial state. 
In Line 7, the operator $Accessible(\cdot)$ extracts the union of the largest component of the BCC that is reachable from the initial state $q^Y_0$. 
It is required that, for each microstate in the returned BCC $\mathcal{C}_c$, the corresponding Y- or Z-state is \emph{consistent}, i.e., all microstates associated with a given information state are included, so that the structural integrity of the BCC is preserved.  
%
For the remainder of the BCC $\mathcal{C}_c$, if there exists at least one Y-microstate corresponding to a cognitively marked state in $\bm{\mathcal{Q}}_{m}$, then $\mathcal{C}_c$ represents a cognitive non-blocking supervisor that is maximally permissive with respect to the marked states.  
Otherwise, no cognitive non-blocking supervisor exists, and the algorithm correctly returns ``No solution''.


\begin{algorithm}[t!]
\caption{Synthesis of Cognitive Non-Blocking BCC}
\label{alg:cnb_syn}
\KwIn{BCC $\mathcal{C} = (\bm{\mathcal{Q}}, \Sigma_\mathcal{C}, \delta_\mathcal{C}, q^Y_0, \bm{\mathcal{Q}}_m)$}
\KwOut{Cognitive non-blocking BCC $\mathcal{C}_c$}

\BlankLine

% \tcp{Step 1: Prune infeasible Z-states}
Identify all Z-microstates in $\bm{\mathcal{Q}}^Z_l \subseteq \bm{\mathcal{Q}}$ that are deadlock or livelock using \eqref{eq:bcc_dl} and \eqref{eq:bcc_ll}\;
Prune the corresponding Z-states $z$, $(z, q) \in \bm{\mathcal{Q}}^Z_l$ from $\bm{\mathcal{Q}}$\;

\Do{BCC $\mathcal{C}$ not consistent}{

    Prune the Y-states that are not cognitively live or violate Definition~\ref{def:con_bcc} (i)\;
    Prune the Z-states that violate Definition~\ref{def:con_bcc} (ii)\;
    
}
$\mathcal{C}_c = (\bm{\mathcal{Q}}_c, \Sigma_\mathcal{C}, \delta_\mathcal{C, c}, q^Y_0, \bm{\mathcal{Q}}_{m,c}) \gets Accessible(\mathcal{C})$\;
\eIf{$\bm{\mathcal{Q}}_c \cap \bm{\mathcal{Q}}_m = \emptyset$}{
    \Return $\mathrm{No \, solution}$\;
}{
    \Return BCC $\mathcal{C}_c$\;
}


\Return BCC $\mathcal{C}$\;
\end{algorithm}


\begin{example}\label{ex:cnb_removal}
    Let us consider the DFA in Fig.~\ref{fig:DFA_aic_overall} (i) and the corresponding BCC in Fig.~\ref{fig:bcc_overall} again.
    We use Algorithm~\ref{alg:cnb_syn} to prune the states that do not satisfy the property of potentially feasible, or cognitive live, and then we make this system consistent.
    %
    First, following the result of Example~\ref{ex:nb_1}, one can easily check that the blocking Z-states $z_2$, $z_3$, $z_6$, $z_7$, $z_8$, and $z_{12}$ do not satisfy the definition of potential feasibility, and should be pruned.
    %
    We also prune $z_{10}$ for its microstate $(z_{10}, q_19)$ is a deadlock state.
    %
    Next, one can easily check that $y_5$ violates the consistency of BCC and should be pruned.
    We also iteratively prune $z_7$ and $z_8$ for the requirement of consistency.
    %
    For Y-states $y_1$ and $y_3$, which are not cognitive live, we need to prune these states, and then we prune $z_4$ for consistency.
    Finally, based on Definition~\ref{def:con_bcc} (ii), all enabled observations must be defined for Z-states, so the Z-state $z_1$ should be removed.
    %
    The removed states are marked by red crosses in Fig.~\ref{fig:bcc_overall}, and the remaining structure forms the resulting consistent cognitive non-blocking BCC.
\end{example}


In the remainder of this article, to simplify the expression, we use $\mathcal{C}_c$ instead of $\mathcal{C}$ to denote the consistent cognitive non-blocking BCC output by Algorithm~\ref{alg:cnb_syn}.

% TODO
% 因为这么删确实会把原系统弄得很破碎
\begin{remark}
Although the resulting BCC output by Algorithm~\ref{alg:cnb_syn} may appear ``fragmented'' in terms of state space, this is inevitable: any state that cannot participate in a feasible trajectory or reach a marked Y-state under partial observation cannot be part of a cognitive non-blocking supervisor. 
To ensure correctness, we consider only the accessible component of the BCC, i.e., states reachable from the initial state $q^Y_0$. 
Thus, pruning such states is both necessary and sufficient to guarantee that the synthesized supervisor enforces a cognitive non-blocking sublanguage.
\end{remark}


\begin{figure*}
    \centering
    \includegraphics[width=0.765\linewidth]{Figs/conceptual_proof.pdf}
    \caption{Conceptual illustration of Proof of Theorem~\ref{theo:cnb}.}
    \label{fig:conceptual_proof}
\end{figure*}

\subsection{Synthesized supervisor and correctness}

In this subsection, we show how the cognitive non-blocking BCC 
$\mathcal{C}$ induces a supervisor and proves the correctness of our synthesis
algorithm.
%
Based on Theorem~\ref{theo:ctrl_bcc}, the projected language of the BCC 
satisfies \eqref{eq:bcc_controllability}, which means that it is closed under
all plant-feasible uncontrollable events with respect to $G$ and
$\Sigma_{uc}$.
%
Therefore, the pruned BCC does not represent an idealized behavior set that can
be followed only when all events are controllable. Instead, it defines a
behavioral structure that can be implemented by a supervisor, since the plant
cannot be forced to leave $\mathcal{L}_{\Sigma}(\mathcal{C})$ through the
occurrence of an uncontrollable event.
%
Moreover, by Algorithm~\ref{alg:cnb_syn}, the remaining structure of
$\mathcal{C}$ preserves all feasible microstate-level paths required for
cognitive non-blocking. Hence, every control-observation evolution encoded in
$\mathcal{C}$ can be interpreted as an admissible evolution of the plant
under supervisory control.

Recall that the projection from the BCC language to the sequence of control patterns is
$
P_{\Gamma}:\Sigma_{\mathcal{C}}^* \rightarrow \Gamma^*,
$
and the set of event sequences generated by BCC is denoted as $\mathcal{L}_{\Sigma}(\mathcal{C})$.
For a BCC run
$
\rho_\mathcal{C}
=
(\bm{q}_0,q_0)(\bm{q}_1,q_1)\cdots(\bm{q}_n,q_n)\cdots,
$
the corresponding plant-event sequence and control-pattern sequence are,
respectively,
$
    P_\Sigma(tr_\mathcal{C}(\rho_\mathcal{C}))\in\Sigma^*,
$
and
$
    P_\Gamma(tr_\mathcal{C}(\rho_\mathcal{C}))\in\Gamma^*,
$
respectively.
Note that the supervisor $S: P(\mathcal{L}(G)) \rightarrow 2^\Gamma$ can be equivalently expressed as the mapping from projected plant sequences to admissible control sequences in BCC, $S_{\mathcal{C}} : P(\mathcal{L}_{\Sigma}(\mathcal{C})) \rightarrow P_\Gamma(\mathcal{L}(\mathcal{C}))$.
%
Consequently, the supervisor $S$ induced by the BCC can be described as a mapping from observation sequences to control-pattern sequences. Specifically, for each observation sequence 
$\alpha \in P(\mathcal{L}_{\Sigma}(\mathcal{C}))$, for a slight abuse of notation, we define
\begin{flalign}\label{eq:bcc_sc:1}
%
% 试过了，这个就是一行放不下
% S_{\mathcal{C}}(\alpha) = \{ &\nu \in \Gamma^* \mid\ \exists \eta \in \mathcal{L}(\mathcal{C}) : \notag \\
% &P(P_\Sigma(\eta)) = \alpha, \nu = P_\Gamma(\eta) \}.\\
%
S_{\mathcal C}(\alpha) :=\{&\gamma\in\Gamma\mid \exists \eta \in \mathcal L(\mathcal C): P(P_\Sigma(\eta))=\alpha, \notag \\
&\eta\gamma\in\mathcal L(\mathcal C)\}.
\end{flalign}
%
%
That is, $S_{\mathcal{C}}(\alpha)$ collects all control-pattern sequences
encoded by BCC runs whose projected plant-event sequences generate the same
observation $\alpha$.
%
The closed-loop language induced by $S_{\mathcal{C}}$ is then defined by the
plant-event projections of the same BCC runs
\begin{flalign}\label{eq:bcc_sc:2}
\mathcal{L}(S_{\mathcal{C}}/G)
:=&
\{\, s\in\Sigma^* \mid
\exists \eta\in\mathcal{L}(\mathcal C) :  s=P_\Sigma(\eta),  \notag\\
& \; \; P_\Gamma(\eta) \in S_{\mathcal{C}}(P(s))
\,\} \notag \\
=&
P_\Sigma(\mathcal{L}(\mathcal C))
=
\mathcal{L}_{\Sigma}(\mathcal C).
\end{flalign}
%
By \eqref{eq:bcc_sc:1}, the condition $P_\Gamma(\eta)\in S_{\mathcal{C}}(P(s))$ means that the control-pattern
projection of the same BCC run $\eta$ is admissible under the observation $P(s)$, while $s=P_\Sigma(\eta)$ gives the corresponding plant-event
projection.
Thus, \eqref{eq:bcc_sc:1} extracts the control-sequence interpretation of a BCC run, whereas \eqref{eq:bcc_sc:2} extracts its event-sequence interpretation.
%
By Proposition~\ref{prop:bcc_dfa_consistency}, every such projected event sequence is a feasible execution of the DFA $G$; moreover, by
Theorem~\ref{theo:ctrl_bcc}, this projected language is closed under all plant-feasible uncontrollable events.
Hence, we can apply $S_\mathcal{C}$ as $S$ such that $\mathcal{L}(S_{\mathcal{C}}/G)$ can be used as an equivalent closed-loop representation of the controlled DES, $\mathcal{L}(S/G)$, without explicitly constructing the supervisor-plant product automaton $S_{\mathcal{C}}\times G$.

Now, we give the following theorem to show that the supervisor synthesized from the BCC guarantees cognitive non-blocking.

\begin{theorem}\label{theo:cnb}
Given a plant $G = (Q, \Sigma, \delta, q_0, Q_m)$ and the cognitive non-blocking BCC $\mathcal{C}$ constructed by Algorithm~\ref{alg:cnb_syn}, 
the supervisor $S_{\mathcal{C}}$ synthesized from $\mathcal{C}$ guarantees that the closed-loop system is cognitive non-blocking.
\end{theorem}
%
\begin{proof}
The proof proceeds in two main steps, and the conceptual illustration is given in Fig.~\ref{fig:conceptual_proof}.
%
First, we relate the closed-loop behavior of $S_{\mathcal{C}}/G$ to the pruned BCC $\mathcal{C}$.
By Proposition~\ref{prop:bcc_dfa_consistency}, every BCC run projects to a feasible execution of the plant $G$.
Moreover, by Theorem~\ref{theo:ctrl_bcc}, 
$\mathcal{L}_{\Sigma}(\mathcal{C})$ is closed under all plant-feasible uncontrollable events.
Thus, we obtain
$$
\forall (\bm{q'}, q') = \delta_\mathcal{C}((\bm{q}, q), \bm{s}), \exists q' = \delta(q, s) :
P_\Sigma(\bm{s}) = s.
$$
%s
% 证明主要分为两部分，第一部分是结构
% 为什么BCC可以代替作为受控系统进行讨论，需要一套符号化说明
%
% 1 算法2剩下的BCC仍然是一个初始状态可达的BCC，所以满足\eqref{eq:con_bcc}
% 2 说明\eqref{eq:con_bcc}当\mathcal{L}_{\Sigma}(\mathcal{C})作为规范语言的话，不会因为不可控事件的发生而被迫离开自身。
% 3 对于BCC，其中任何Z-Z transition和Z-Y transition，都可以在自动机中找到对应的结果，且其变迁完全与自动机对应
% 基于上述三点，相当于说明了BCC可以代替作为受控系统进行讨论
%

% 
Then, we prove that the structure of BCC satisfies the notion of cognitive non-blocking.
%
Let $\bm{\mathcal{Q}}_m \subseteq \bm{\mathcal{Q}}^Y$ denote the cognitively marked Y-microstates.  
By Algorithm~\ref{alg:cnb_syn}, $\mathcal{C}$ is constructed such that:
%
\begin{enumerate}[(i)]
    \item All deadlock and livelock Z-microstates have been pruned, ensuring that for any $(z,q) \in \bm{\mathcal{Q}}^Z$ in $\mathcal{C}$, there exists a path to a Y-microstate $(y, q)$ corresponding to the successor Y state $y$;
    \item All Y-states are cognitively live with respect to $\bm{\mathcal{Q}}_m$, i.e., for each $y \in \bm{\mathcal{Q}}^Y$ and all $q \in y$, there exists a marked Y-microstate $(y_m, q_m) \in \bm{\mathcal{Q}}_m$ such that
    $
        (y,q) \leadsto_{\bm{\mathcal{Q}}} (y_m,q_m);
    $  
    \item All microstates are accessible from the initial state $q^Y_0 = (y_0, q_0)$.
\end{enumerate}
%

To this end, given an arbitrary microstate $(\bm{q}, q) \in \bm{\mathcal{Q}}$, there exists the sequence $\bm{s}_{q_0, q}$ in BCC $\mathcal{C}$ such that 
$$
(\bm{q}, q) = \delta_{\mathcal{C}}((y_0, q_0), \bm{s}_{q_0, q}),
$$
and the corresponding sequence of events can be denoted as $P_\Sigma(\bm{s}_{q_0, q}) = s$.
%
Similarly, with the co-accessibility in BCC guaranteed, we have $\bm{t}_{q, q_m}$ such that $\mathcal{C}$ such that 
$$
(y_m, q_m) = \delta_{\mathcal{C}}((\bm{q}, q), \bm{t}_{q, q_m}),
$$
and the corresponding sequence of events can be denoted as $P_\Sigma(\bm{t}_{q, q_m}) = t$.
%
Based on Definition~\ref{def:con_bcc}, for state $(\bm{q}, q)$, if it is a Z-microstate, then it can be guaranteed to reach a Y-microstate of a successor Y-state.
% 
Based on \eqref{eq:bcc_y_cl}, given a cognitively live Y-state, all its Y-microstates can reach a cognitively marked state via the sequences corresponding to identical Y- and Z-states.
%
Moreover, based on the bipartite structure inherited from the AIC, for any microstate $(\bm{q}, q')$ that shares identical observations with state $(\bm{q}, q)$, such that for
$$
\begin{aligned}
    (\bm{q}, q') &= \delta_{\mathcal{C}}((y_0, q_0),   \bm{s}_{q_0, q'}) \\
    (y_m,   q_m) &= \delta_{\mathcal{C}}((\bm{q}, q'), \bm{t}_{q', q_m})
\end{aligned}
$$
we have
$
P(P_\Sigma(\bm{s}_{q_0, q'})) = P(P_\Sigma(\bm{s}_{q_0, q}))
$
and
$
P(P_\Sigma(\bm{t}_{q', q_m})) = P(P_\Sigma(\bm{t}_{q, q_m})).
$
Formally, for any $s \in \mathcal{L}_{\Sigma}(\mathcal{C})$, there exists a continuation $t \in \Sigma^*$ such that the corresponding cognitive state estimate satisfies
$
\hat{\mathcal{E}}_{S_{\mathcal{C}}}(st) \subseteq \bm{\mathcal{Q}}_m.
$
%
Thus, in the corresponding DFA $G$ for any $s \in \mathcal{L}(S_{\mathcal{C}} / G)$, there exists a continuation $t \in \Sigma^*$ such that the corresponding cognitive state estimate satisfies
$$
\hat{\mathcal{E}}_{S_{\mathcal{C}}}(st) \subseteq Q_m.
$$
This completes the proof.
\end{proof}

\begin{figure}[ht!]
    \centering
    \includegraphics[width=0.65\linewidth]{Figs/controlled_fsa_overall.pdf}
    \caption{Controlled DFA $S / G$.}    % TODO 这个符号对不对
    \label{fig:ctrl_DFA}
\end{figure}

We give Example~\ref{ex:DFA_cnb} to illustrate that the closed-loop system satisfies the notion of cognitive non-blocking.
Meanwhile, we discuss the complexity of Algorithm~\ref{alg:cnb_syn}.

\begin{example}\label{ex:DFA_cnb}
    Let us consider the DFA illustrated in Fig.~\ref{fig:DFA_aic_overall} again,  and the corresponding BCC illustrated in Fig.~\ref{fig:bcc_overall}.
    %
    After pruning the inconsistent states, the corresponding closed-loop system $S / G$ is illustrated in Fig.~\ref{fig:ctrl_DFA}.
    %
    In the resulting closed-loop system, the evolutions from the initial state $q_0$ are $q_0 \xrightarrow{u_1} q_1 \xrightarrow{o_2} q_5$ and $q_0 \xrightarrow{u_2} q_2 \xrightarrow{o_2} q_6$, and the corresponding observation is $o_2$.
    In the cognitive sense, after observing $o_2$, the supervisor can certify that the current plant state belongs to the marked set
    \begin{equation}\label{eq:ex_cnb_reach_marked}
        \hat{\mathcal{E}}_S(o_2) = \{ q_5, q_6 \} \subseteq Q_m,
    \end{equation}
    although the exact state may still be unknown.
    %
    Moreover, from any reachable state after $q_5$ or $q_6$, the controlled system can evolve through a finite admissible continuation and return to $q_0$.
    The particular observable labels along this return segment are not essential for the cognitive non-blocking argument.
    After returning to $q_0$, the continuation in \eqref{eq:ex_cnb_reach_marked} can be applied again to reach the cognitively marked information state $\{ q_5, q_6 \}$.
    %
    Hence, for any sequence in $G$ corresponding to a cognitively live path in the BCC $\mathcal{C}$, there exists a continuation that reaches a cognitively marked state.
    Once the supervisor $S$ is induced from the BCC $\mathcal{C}$ as described in Section~\ref{subsec:supervisor_induction}, the closed-loop system $S/G$ inherits this property and satisfies cognitive non-blocking.
    
    Let us consider the motivated example again in Section~\ref{sec:intro}.
    The robotic system cannot certify its exact current state, but its state estimate can be updated after completing certain motion primitives.
    When it is at the initial state $q_0$, it should reject the motions corresponding to $o_3$, which may lead to $q_3$ or $q_4$.
    In this sense, it can then arrive at $q_5$ or $q_6$; moreover, when it observes the information corresponding to event $o_2$, it knows that it is in a safe location, so the transmission of critical information can be started.
\end{example}

\begin{proposition}
Consider a BCC $\mathcal{C}$ constructed over a plant with $|Q|$ states and
event set $\Sigma$.  the worst-case time complexity of Algorithm~\ref{alg:cnb_syn} is
$
    O\big(|Q|\cdot |\Sigma|\cdot 2^{2|Q|+2|\Sigma|}\big).
$
\end{proposition}

\begin{proof}
Algorithm~\ref{alg:cnb_syn} proceeds by pruning Y- and Z-states of the BCC.
In Lines 1--2, deadlock and livelock Z-microstates are identified by a BFS
search, which can be done in time
$
O(|\bm{\mathcal{Q}}|+|\delta_{\mathcal{C}}|)
$.
%
In Lines 3--6, each iteration prunes Y-states that are not cognitively live
or Z-states that violate consistency. Each iteration can be implemented by one
BFS over the remaining BCC, with cost
$
O(|\bm{\mathcal{Q}}|+|\delta_{\mathcal{C}}|)
$.
Since each nontrivial iteration prunes at least one Y-state or Z-state, the
number of iterations is at most $|Y|+|Z|$. Therefore, the total cost is
$
    O((|Y|+|Z|)(|\bm{\mathcal{Q}}|+|\delta_{\mathcal{C}}|)).
$
Line 7 performs one additional BFS, which is dominated by the above term.
Hence, the overall complexity is
$
    O\big((|Y|+|Z|)(|\bm{\mathcal{Q}}|+|\delta_{\mathcal{C}}|)\big).
$
%
By Proposition~\ref{prop:state_bcc}, we have
$
    |\bm{\mathcal{Q}}|
    \le
    |Q|\big(2^{|Q|}+2^{|Q|+|\Sigma|}\big)
    =
    O(|Q|\cdot 2^{|Q|+|\Sigma|}),
$
and
$$
    |\delta_{\mathcal{C}}|
    \le
    |Q|2^{|Q|+|\Sigma|}(1+|\Sigma|)
    =
    O(|Q|\cdot |\Sigma|\cdot 2^{|Q|+|\Sigma|}).
$$
Therefore, based on
$
|Y|+|Z| = O(2^{|Q|+|\Sigma|}),
$
we have
$$
\begin{aligned}
    O\big((|Y|+|Z|)(|\bm{\mathcal{Q}}|+|\delta_{\mathcal{C}}|)\big)
    &=
    O\big(
    2^{|Q|+|\Sigma|}
    \cdot
    |Q|\cdot |\Sigma|\cdot 2^{|Q|+|\Sigma|}
    \big) \\
    &=
    O\big(|Q|\cdot |\Sigma|\cdot 2^{2|Q|+2|\Sigma|}\big).
\end{aligned}
$$
This completes the proof.
\end{proof}
%


\begin{corollary}
    The cognitive non-blocking synthesis problem is solvable by Algorithms~\ref{alg:bcc} and \ref{alg:cnb_syn} within finite steps.
\end{corollary}



% \section{Case study}

% \begin{figure}
%     \centering
%     \includegraphics[width=\linewidth]{Figs/fsa_case.pdf}
%     \caption{DFA $G_1$ of  case studies, with observable event set $\Sigma_o = \{ o_1, o_2 \}$ and controllable events $\Sigma_o = \{ c_1, c_2 \}$.}
%     \label{fig:DFA_case}
% \end{figure}

% \begin{figure}
%     \centering
%     \includegraphics[width=\linewidth]{Figs/aic_case.pdf}
%     \caption{AIC $\mathcal{A}_1$ constructed from DFA $G_1$ of Fig. \ref{fig:DFA_case}.}
%     \label{fig:aic_case}
% \end{figure}

% \begin{figure}
%     \centering
%     \includegraphics[width=0.785\linewidth]{Figs/bcc_case.pdf}
%     \caption{BCC $\mathcal{C}_1$ constructed from DFA $G_1$ and AIC $\mathcal{A}_1$ of Figs. \ref{fig:DFA_case} and Figs. \ref{fig:aic_case}.}
%     \label{fig:bcc_case}
% \end{figure}

\section{Special case: marked tasks within the unobservable reaches}
\label{sec:special_z_marked}

%
% 这个系统要处理的话
% 其实方法是一样的
% 1 对于marked Z-state，需要当前Z-状态中的所有状态能够到达这个状态，或者退到下一个Y状态
% 2 Y状态不变
% 3 系统要consistent
% 

In this section, we discuss a special case of the proposed method, namely, systems in which marked tasks may be completed within unobservable reaches.
In the previous sections, the cognitively marked states are defined over Y-microstates, since Y-states represent information states immediately after observation updates.
However, in some systems, a marked state can be reached inside a Z-state before the next observable event occurs.
%
The main result of this article already handles the case where task completion is certified at marked Y-microstates. Therefore, we only briefly discuss the case of marked Z-microstates here.

To handle this case, we extend the target set of the BCC by including marked Z-microstates. 
For a Z-state $z$, define
$
    \bm{\mathcal{Q}}_{z,m} = \{(z,q)\in \bm{\mathcal{Q}}_z \mid q\in Q_m\}
$
as the set of marked Z-microstates in $z$. 
Meanwhile, the set of exit Z-microstates can be defined as
$$
    \bm{\mathcal{Q}}_{z,e}
    =
    \{ (z, q_e) \in \bm{\mathcal{Q}}_z \mid
    \exists \sigma_o \in (\Gamma(z) \cup \Sigma_{uc}) \cap \Sigma_o : \delta(q_e,\sigma_o)! \}.
$$
%
We say that a Z-state $z$ is \emph{a marked Z-state} if (i) for every microstate $(z,q)\in \bm{\mathcal{Q}}_z$, there exists a marked Z-microstate $(z,q_m)\in \bm{\mathcal{Q}}_{z,m}$ such that
$$
    (z,q) \leadsto_{\bm{\mathcal{Q}}_z} (z,q_m).
$$
(ii) for this marked Z-microstate $(z, q_m)$, there exists an exit Z-microstate $(z, q_e)\in \bm{\mathcal{Q}}_{z,e}$ and a path inside $\bm{\mathcal{Q}}_z$
$$
    \rho_{\mathcal{C}}
    =
    (z,q_0)(z,q_1)\cdots(z,q_\ell),
    \quad
    q_0=q_m, \, q_\ell=q_e,
$$
such that, for each $i=0,\ldots,\ell-1$, there exists $\sigma_{uo,i} \in \Sigma_{uo}$ satisfying
$
    \delta_{\mathcal{C}}((z,q_i),\sigma_{uo,i}) = (z,q_{i+1}),
$
and
$
    Post((z,q_i))\cap \bm{\mathcal{Q}}_z = \{ (z, q_{i+1}) \}.
$
Hence, after reaching the marked Z-microstate $(z,q_m)$, the subsequent internal evolution inside $z$ is uniquely determined until the exit Z-microstate $(z,q_e)$ is reached.
%
In other words, each microstate of a marked Z-state can reach a marked Z-microstate, after which the internal path to the exit is unique. 
Thus, the observed exit can certify task completion within the unobservable reach.

We use $Z_m$ to denote the set of marked Z-states.
Then, the target set for the cognitive-liveness is extended from
$\bm{\mathcal{Q}}_m$ to
\begin{equation}
    \bm{\mathcal{Q}}_m^{+}
    =
    \bm{\mathcal{Q}}_m
    \cup
    \bigcup_{z\in Z_m}\bm{\mathcal{Q}}_{z,m}.
    \label{eq:extended_target_set}
\end{equation}
Accordingly, the definition of cognitively live Y-states in
\eqref{eq:bcc_y_cl} is extended as
\begin{equation}
\label{eq:bcc_y_cl:2}
    \forall q \in y,\;
    \exists (\bm q_t,q_t)\in \bm{\mathcal{Q}}_m^{+}
    :
    (\bm q_t,q_t)
    =
    \delta_{\mathcal{C}}
    \big((y,q),tr(\rho_{\mathcal{C}}')\big),
\end{equation}
where the same BCC sequence $tr(\rho_{\mathcal{C}}')$ is used for all
$q\in y$, as in \eqref{eq:bcc_y_cl}.
Thus, a cognitively live Y-state is allowed to reach either a cognitively
marked Y-microstate or a marked Z-microstate whose internal path to the
next observable exit.
%
Based on these modifications, Algorithm~\ref{alg:cnb_syn} can still be
applied to obtain a consistent BCC. The resulting BCC then induces a
cognitive non-blocking supervisor for systems in which task completion can
be certified either at marked Y-microstates or within marked Z-states.
In this way, Problem~\ref{prob:cnb} can be handled by replacing
\eqref{eq:est_o} with \eqref{eq:esti:1}.

It should be noted that completion of a marked Z-state is certified only after the subsequent observable exit is observed. 
This is sufficient for the mathematical definition, since all compatible BCC paths must have visited a marked Z-microstate before the exit. 
However, in many practical applications, the unobservable tail inside the Z-state is not physically observed, and thus the exact occurrence of task completion may not be directly interpretable from the subsequent observation alone.

\section{Conclusion}\label{sec:con}

In this article, we proposed a new non-blocking property, called cognitive non-blocking, under which the observer can certify that the system is able to reach a marked state. 
Based on this notion, we developed a supervisor synthesis algorithm that constructs a controlled sublanguage guaranteed to be cognitive non-blocking. 
We also proved that the proposed synthesis algorithm terminates in finite steps. 
%
Future work will investigate how the notion of cognitive non-blocking can be extended to weighted and interval-weighted systems, where task completion may depend not only on reachability but also on accumulated transition weights.

\section*{References}\vspace{-1em}
\bibliographystyle{IEEEtran}
\bibliography{ref_abbr}


\begin{IEEEbiography}[{\includegraphics[width=1in,height=1.25in,clip,keepaspectratio]{a1.png}}]{First A. Author} (Fellow, IEEE) and all authors may include 
biographies. Biographies are
often not included in conference-related articles.
This author is an IEEE Fellow. The first paragraph
may contain a place and/or date of birth (list
place, then date). Next, the author’s educational
background is listed. The degrees should be listed
with type of degree in what field, which institution,
city, state, and country, and year the degree was
earned. The author’s major field of study should
be lower-cased.

The second paragraph uses the pronoun of the person (he or she) and
not the author’s last name. It lists military and work experience, including
summer and fellowship jobs. Job titles are capitalized. The current job must
have a location; previous positions may be listed without one. Information
concerning previous publications may be included. Try not to list more than
three books or published articles. The format for listing publishers of a book
within the biography is: title of book (publisher name, year) similar to a
reference. Current and previous research interests end the paragraph.

The third paragraph begins with the author’s title and last name (e.g.,
Dr. Smith, Prof. Jones, Mr. Kajor, Ms. Hunter). List any memberships in
professional societies other than the IEEE. Finally, list any awards and work
for IEEE committees and publications. If a photograph is provided, it should
be of good quality, and professional-looking.
\end{IEEEbiography}

\begin{IEEEbiographynophoto}{Second B. Author,} photograph and biography not available at the
time of publication.
\end{IEEEbiographynophoto}

\begin{IEEEbiographynophoto}{Third C. Author Jr.} (Member, IEEE), photograph and biography not available at the
time of publication.
\end{IEEEbiographynophoto}

\end{document}
